\documentclass[12pt,a4paper,UTF8]{ctexart}
\usepackage[a4paper,top=22mm,bottom=22mm,left=22mm,right=22mm,headheight=15pt]{geometry}
\usepackage{amsmath,amssymb,mathtools,bm}
\usepackage{graphicx,xcolor,booktabs,longtable,array,enumitem,fancyhdr,caption,float,microtype,setspace}
\usepackage[hidelinks,unicode]{hyperref}
\usepackage[most]{tcolorbox}
\definecolor{ContestBlue}{HTML}{17365D}
\hypersetup{pdftitle={第086套},pdfauthor={}}
\setstretch{1.24}
\setlength{\parindent}{2em}
\setlength{\parskip}{0.35em}
\setlength{\emergencystretch}{2em}
\setlist{itemsep=0.35em,topsep=0.4em}
\captionsetup{font=small,labelfont=bf,justification=centering}
\setcounter{secnumdepth}{0}
\pagestyle{fancy}
\fancyhf{}
\fancyhead[L]{\small 第086套}
\fancyfoot[C]{\thepage}
\renewcommand{\headrulewidth}{0.4pt}
\ctexset{
  section={format=\Large\bfseries\color{ContestBlue},beforeskip=1.2em,afterskip=0.7em},
  subsection={format=\large\bfseries\color{ContestBlue},beforeskip=1.0em,afterskip=0.5em},
  subsubsection={format=\normalsize\bfseries,beforeskip=0.8em,afterskip=0.35em}
}
\makeatletter
\newcommand{\tightlist}{\setlength{\itemsep}{0pt}\setlength{\parskip}{0pt}}
\newcommand{\maxwidth}{\ifdim\Gin@nat@width>\linewidth\linewidth\else\Gin@nat@width\fi}
\newcommand{\maxheight}{\ifdim\Gin@nat@height>0.72\textheight 0.72\textheight\else\Gin@nat@height\fi}
\newcommand{\pandocbounded}[1]{\begingroup\setkeys{Gin}{width=\maxwidth,height=\maxheight,keepaspectratio}#1\endgroup}
\makeatother
\begin{document}
\begin{titlepage}
  \centering
  \vspace*{58mm}
  {\Huge\bfseries\color{ContestBlue} 第086套\par}
  \vfill
  \vspace*{28mm}
\end{titlepage}
\section{第 1 题}\label{ux7b2c-1-ux9898}

如图，在一个足够长的光滑斜面上连接一个边长为 \(d\)
的正方形光滑水平平台，斜面的倾角为 \(30^\circ\)。水平平台上有一质量为
\(m\)
的小球。轻绳的一端系着小球，穿过平台中心的光滑小孔，另一端系在斜面上一质量为
\(M=2m\) 的小物块上，物块可视为质点。绳长为 \(1.5d\)，重力加速度为
\(g\)。

\begin{enumerate}
\def\labelenumi{\arabic{enumi}.}
\tightlist
\item
  若小球绕小孔做半径为 \(d/4\) 的匀速圆周运动，求角速度 \(\omega_0\)；
\item
  在上述运动状态下，突然给物块一个沿斜面向下的冲量，使物块获得速度
  \(v'=\sqrt{2gd}\)，求此时小球指向圆心的径向速度；
\item
  此后若物块始终与斜面接触，求物块沿斜面向下运动的最大位移，结果给出数值解；
\item
  重新回到第（1）问的状态。沿小球初速度方向再给小球一个冲量
  \(I\)，使其切向速度瞬间增大。求当 \(I\)
  超过多大时，小球有可能从正方形平台边缘掉下。
\end{enumerate}

\begin{figure}[H]
\centering
\includegraphics[width=0.502\linewidth,height=\textheight,keepaspectratio,alt={图 1-1　水平平台、斜面、轻绳与两物体的空间位置。}]{../figures/question-01.pdf}
\caption*{图 1-1　水平平台、斜面、轻绳与两物体的空间位置。}
\end{figure}

\subsection{一、题型定位}\label{ux4e00ux9898ux578bux5b9aux4f4d}

小球的径向运动会带动斜面物块，冲量后的径向速度不能只用小球质量计算。先用直绳几何求速度关系，再用角动量与能量判断转向点及临界冲量。

\subsection{二、知识点}\label{ux4e8cux77e5ux8bc6ux70b9}

\begin{enumerate}
\def\labelenumi{\arabic{enumi}.}
\tightlist
\item
  光滑小孔两侧轻绳张力相等；
\item
  斜面内三角形的正弦定理与勾股关系；
\item
  不可伸长绳的速度约束，即绳长对时间求导；
\item
  关于小孔的角动量守恒；
\item
  重力、动能和几何约束共同决定的转向点；
\item
  冲量只改变切向速度时的角动量与能量关系；
\item
  平台边界对应的临界半径 \(r=d/2\)。
\end{enumerate}

\subsection{三、物理图像}\label{ux4e09ux7269ux7406ux56feux50cf}

平台中心小孔记为 \(O\)，平台与斜面的交界处取距 \(O\) 最近的点
\(A\)，物块位置记为 \(B\)。平台中心到交界边的垂直距离为
\(OA=d/2\)。当小球半径为 \(r\) 时，平台上绳段长为 \(r\)，因总绳长为
\(3d/2\)，悬空直绳 \(OB\) 的长度为 \(3d/2-r\)；绳没有在边缘 \(A\)
处绕过导向轮。

小球所受水平力始终沿 \(O\) 指向径向，因此它关于 \(O\)
的力矩为零；物块的运动改变斜面内绳段长度，进而迫使小球产生径向速度。小球的切向运动与径向运动因此具有不同的守恒规律。

\subsection{四、解题入口}\label{ux56dbux89e3ux9898ux5165ux53e3}

第（1）问先由固定几何求绳在斜面内与斜面方向的夹角
\(\beta\)，再对物块沿斜面列平衡方程。第（2）问只考察冲量后的速度瞬间，对绳长约束求导即可。第（3）问把物块最大下移位置视为径向转向点，用几何方程、角动量守恒和机械能守恒联立。第（4）问先确定小球最先可能越过的边界半径，再由能量求临界切向速度，最后用冲量---动量定理减去原有速度。

\subsection{五、解析}\label{ux4e94ux89e3ux6790}

图 1-2 上方画出直绳 \(OB\) 与斜面投影的直角三角形，\(H\)
是小孔到斜面的垂足；下方俯视图分开画出小球的切向速度与指向小孔的径向速度。先由
\(OH=d/4\) 与绳长求 \(\beta\)，再使用角动量和能量计算转向点与平台边界。

\begin{figure}[H]
\centering
\includegraphics[width=0.502\linewidth,height=\textheight,keepaspectratio,alt={图 1-2　斜面绳长约束与小球的径向、切向速度。}]{../figures/solution-01.pdf}
\caption*{图 1-2　斜面绳长约束与小球的径向、切向速度。}
\end{figure}

\textbf{（1）求稳态角速度}

在初始稳态中，小球半径为

\[
r_0=\frac d4.
\]

平台上绳段长度为 \(d/4\)，所以斜面一侧所对应的总长度为

\[
\frac{3d}{2}-\frac d4=\frac{5d}{4}.
\]

在斜面内考察三角形 \(OAB\)。\(OA=d/2\)，而 \(OB=5d/4\)。设 \(OB\)
与斜面方向的夹角为 \(\beta\)。由正弦定理，

\[
\frac{\sin\beta}{d/2}
=\frac{\sin30^\circ}{5d/4}.
\]

逐项化简：

\[
\begin{aligned}
\sin\beta
&=\frac{d}{2}\cdot\frac{\sin30^\circ}{5d/4}\\
&=\frac d2\cdot\frac{1/2}{5d/4}\\
&=\frac15.
\end{aligned}
\]

因此

\[
\cos\beta
=\sqrt{1-\sin^2\beta}
=\sqrt{1-\frac1{25}}
=\frac{2\sqrt6}{5}.
\tag{1}
\]

对小球沿径向列牛顿第二定律。绳张力 \(T\) 提供向心力：

\[
T=m\omega_0^2r_0
=m\omega_0^2\frac d4.
\tag{2}
\]

物块在斜面上静止。沿斜面向下的重力分量为

\[
Mg\sin30^\circ
=2mg\cdot\frac12
=mg.
\]

张力沿斜面向上的分量为 \(T\cos\beta\)，故

\[
T\cos\beta=mg.
\tag{3}
\]

把式（1）、式（2）代入式（3）：

\[
m\omega_0^2\frac d4\cdot\frac{2\sqrt6}{5}=mg.
\]

约去 \(m\)，再解出 \(\omega_0^2\)：

\[
\omega_0^2
=\frac{4g}{d}\cdot\frac5{2\sqrt6}
=\frac{10g}{\sqrt6\,d}.
\]

所以

\[
\boxed{\omega_0=\sqrt{\frac{10g}{\sqrt6\,d}}}.
\tag{4}
\]

\textbf{（2）冲量后小球的径向速度}

冲量作用时间极短，物块位置和绳的几何方向来不及改变，因此仍使用
\(\beta\)。不可伸长绳要求斜面一端沿绳方向缩短多快，平台一端就必须以同样速率缩短。

物块速度沿斜面向下，其在绳方向上的分量为

\[
v'\cos\beta.
\]

这等于小球指向 \(O\) 的径向速度 \(v_r\)，故

\[
\begin{aligned}
v_r
&=v'\cos\beta\\
&=\sqrt{2gd}\cdot\frac{2\sqrt6}{5}\\
&=\frac{2\sqrt{12}}5\sqrt{gd}\\
&=\frac45\sqrt{3gd}.
\end{aligned}
\]

因此

\[
\boxed{v_r=\frac45\sqrt{3gd}},
\]

方向指向小孔 \(O\)。

\textbf{（3）物块沿斜面向下的最大位移}

设物块沿斜面向下移动 \(\delta s\)，小球到小孔的距离由 \(r_0\) 变为
\(r\)。初始时物块到交界点的沿斜面投影长度为

\[
\begin{aligned}
s_0
&=\left(\frac{3d}{2}-\frac d4\right)\cos\beta\\
&=\frac{5d}{4}\cdot\frac{2\sqrt6}{5}\\
&=\frac{\sqrt6}{2}d.
\end{aligned}
\tag{5}
\]

物块下移后，该投影变为 \(s_0+\delta s\)。由斜面内直角三角形，

\[
(s_0+\delta s)^2+\left(\frac d4\right)^2
=\left(\frac{3d}{2}-r\right)^2.
\tag{6}
\]

因此半径 \(r\) 可写成

\[
r=\frac{3d}{2}
-\sqrt{(s_0+\delta s)^2+\frac{d^2}{16}}.
\tag{7}
\]

小球所受张力始终通过 \(O\)，关于 \(O\) 的角动量守恒。初始切向速度为
\(\omega_0r_0\)，转向点切向速度为 \(\omega r\)，所以

\[
m\omega_0r_0^2=m\omega r^2.
\tag{8}
\]

从式（8）解出转向点的切向速度：

\[
\omega r=\frac{\omega_0r_0^2}{r}.
\tag{9}
\]

冲量刚结束时，小球有切向速度 \(\omega_0r_0\) 和径向速度 \(v_r\)；质量为
\(2m\) 的物块速度为
\(v'=\sqrt{2gd}\)。取冲量后的位置为重力势能零点，初始总动能为

\[
K_i
=\frac12m\left[(\omega_0r_0)^2+v_r^2\right]
+\frac12(2m)v'^2.
\tag{10}
\]

物块运动到最低点时，物块速度为零，由绳约束小球径向速度也为零，只剩切向速度。物块下降的竖直高度为
\(\delta s\sin30^\circ\)，系统末态机械能为

\[
E_f
=\frac12m(\omega r)^2
-2mg\,\delta s\sin30^\circ.
\tag{11}
\]

机械能守恒 \(K_i=E_f\)。代入式（9）并用 \(\sin30^\circ=1/2\)：

\[
\frac12m\left[(\omega_0r_0)^2+v_r^2\right]
+mv'^2
=\frac12m\left(\frac{\omega_0r_0^2}{r}\right)^2
-mg\delta s.
\tag{12}
\]

把 \(r_0=d/4\)、式（4）、\(v_r=\frac45\sqrt{3gd}\)、\(v'=\sqrt{2gd}\)
以及式（7）全部代入。为便于数值求解，令

\[
x=\frac{\delta s}{d},
\qquad
\rho=\frac rd.
\]

则

\[
\rho=\frac32-\sqrt{\left(\frac{\sqrt6}{2}+x\right)^2+\frac1{16}}.
\tag{13}
\]

式（12）除以 \(mgd\) 后成为无量纲方程

\[
\frac12\left[
\frac{1}{16}\frac{10}{\sqrt6}
+\frac{48}{25}
\right]+2
=\frac12
\left[
\frac{\frac14\sqrt{10/\sqrt6}\cdot\frac14}{\rho}
\right]^2-x.
\tag{14}
\]

等价地，

\[
\frac{5}{16\sqrt6}+\frac{24}{25}+2
=\frac{5}{256\sqrt6\,\rho^2}-x.
\tag{15}
\]

联立式（13）、式（15）求物理区间 \(0<x<0.254\) 内的根。用二分法计算：

\[
\begin{array}{c|c}
x & \text{式（15）右端减左端}\\ \hline
0.2040 & -4.36\times10^{-2}\\
0.2045 & 2.14\times10^{-2}
\end{array}
\]

继续缩小区间可得

\[
x\simeq0.2043.
\]

按题目要求保留三位有效数字，

\[
\boxed{\delta s_{\max}\simeq0.204d}.
\tag{16}
\]

\textbf{（4）小球越过平台边缘的临界冲量}

小孔位于正方形中心，小球从中心出发最先可能到达的边界是任意一条边，其中点到中心的距离为

\[
r_f=\frac d2.
\tag{17}
\]

初态斜面中物块到交界点的长度为

\[
s_0=\sqrt{\left(\frac{5d}{4}\right)^2-\left(\frac d4\right)^2}
=\frac{\sqrt6}{2}d.
\tag{18}
\]

临界末态平台上绳段长为 \(d/2\)，故斜面一侧总绳长为
\(d\)，相应沿斜面长度为

\[
s_f=\sqrt{d^2-\left(\frac d4\right)^2}
=\frac{\sqrt{15}}4d.
\tag{19}
\]

小球向外移动时，物块沿斜面上升，竖直上升高度

\[
\begin{aligned}
h
&=(s_0-s_f)\sin30^\circ\\
&=\frac12\left(\frac{\sqrt6}{2}-\frac{\sqrt{15}}4\right)d\\
&=\frac{2\sqrt6-\sqrt{15}}8d.
\end{aligned}
\tag{20}
\]

设冲量后小球的角速度为 \(\omega_0'\)，到达边界时为
\(\omega_f\)。角动量守恒给出

\[
m\omega_0'r_0^2=m\omega_fr_f^2.
\tag{21}
\]

由于 \(r_f=2r_0\)，

\[
\omega_fr_f
=\omega_0'\frac{r_0^2}{r_f}
=\frac12\omega_0'r_0.
\tag{22}
\]

临界状态中，小球恰能到达 \(r_f\)，物块在该处的径向速度为零。机械能守恒：

\[
\frac12m(\omega_0'r_0)^2
=\frac12m(\omega_fr_f)^2+2mgh.
\tag{23}
\]

代入式（22）：

\[
\frac12m(\omega_0'r_0)^2
=\frac18m(\omega_0'r_0)^2+2mgh.
\]

移项：

\[
\frac38m(\omega_0'r_0)^2=2mgh,
\]

所以冲量后的临界切向速度为

\[
\omega_0'r_0=4\sqrt{\frac{gh}{3}}.
\tag{24}
\]

冲量沿原切向，因而

\[
I=m\left(\omega_0'r_0-\omega_0r_0\right).
\tag{25}
\]

将式（4）、式（20）、\(r_0=d/4\) 代入：

\[
\frac{I}{m\sqrt{gd}}
=\frac4{\sqrt3}
\sqrt{\frac{2\sqrt6-\sqrt{15}}8}
-\frac14\sqrt{\frac{10}{\sqrt6}}.
\tag{26}
\]

数值为

\[
\boxed{I_{\rm c}\simeq0.3219\,m\sqrt{gd}}.
\tag{27}
\]

因此当

\[
\boxed{I>I_{\rm c}}
\]

时，小球在适当的径向相位下有可能到达并越过平台边缘。若中间几何量过早舍入，临界值会出现约
\(1\%\) 的数值漂移；保留式（26）中的未舍入根式逐项计算，得到式（27）。

\textbf{径向约束的有效质量与数值根检查}

用 \(s\) 表示物块沿坡位移，\(OB^2=s^2+(d/4)^2\)，而
\(r=3d/2-OB\)。微分给
\(\dot r=-\dot s\cos\beta\)。径向冲量不能只改变小球：物块也由不可伸长的绳立刻获得速度。径向动能因此为

\[K_r=\frac12\left(m+\frac{M}{\cos^2\beta}\right)\dot r^2.\]

若对小球施加纯径向冲量
\(J_r\)，\(\Delta\dot r=J_r/(m+M/\cos^2\beta_0)\)。这解释了变式中的有效质量，也说明正文的有效势只描述势能与切向动能；完整径向动力学还要保留位置相关的质量系数。

求第3问数值根时，应同时回代绳长、角动量和能量三式。\(r>0\)
给出允许位移上限。势能随物块下移降低，而切向动能随 \(1/r^2\)
增加；到达转向点时 \(\dot r=\dot s=0\)，所以并不意味着小球瞬时静止。

\textbf{把转向点方程写成一条可检查的函数}

取 \(x=\delta s/d\)，几何半径
\(\rho(x)=3/2-\sqrt{(\sqrt6/2+x)^2+1/16}\)。从初始的 \(\rho(0)=1/4\)
出发，\(\rho\) 随 \(x\) 减小，不能在方程中将它保持为常数。把能量式整理成

\[f(x)=\frac{5}{256\sqrt6\,\rho(x)^2}-x-\left(\frac{5}{16\sqrt6}+\frac{24}{25}+2\right).\]

初始径向冲量后的能量由括号固定，而前两项是转向点的切向动能与重力势。\(f(0)<0\)，当
\(\rho\to0^+\) 时 \(f\to+\infty\)，故至少有一个转向点。数值根约
\(x=0.2043\)，回代可得
\(\rho\approx0.0493\)；小球半径明显缩小，所以角速度显著增大，但角动量保持不变。

对第4问，则先指定临界末态
\(r=d/2\)，反算使径向速度恰为零的初始切向速度。这两问一个未知末位置、一个未知初始动量，不能把同一个数值根直接复用。每次均先从几何求物块高度差，才把它放入能量式。若物块已经越过斜面上绳可直连的范围，原几何将失效，应另检验题设允许位置。

\textbf{切向临界冲量的完整回代}

设平台边界为 \(r_b=d/2\)，初始半径仍为
\(r_0=d/4\)。冲量只沿切向，瞬间位置不变，径向速度初值为零。初始物块位置的坡向投影是
\(s_0=\sqrt6d/2\)；边界绳长 \(OB=d\)，因此末投影为
\(s_b=\sqrt{d^2-(d/4)^2}=\sqrt{15}d/4\)。物块沿坡上升量
\(s_0-s_b\)，高度上升一半：

\[\Delta h=\frac12(s_0-s_b).\]

小球角动量给
\(v_{\theta,b}=v_{\theta,0}'r_0/r_b=v_{\theta,0}'/2\)。临界时末径向速度为零，能量成为

\[\frac12m(v_{\theta,0}')^2=\frac12m(v_{\theta,0}'/2)^2+2mg\Delta h.\]

整理得 \((v_{\theta,0}')^2=(8g/3)(s_0-s_b)\)。原始切向速度为
\(v_{\theta,0}=\omega_0d/4\)，故

\[J_c=m(v_{\theta,0}'-v_{\theta,0})\simeq0.3219m\sqrt{gd}.\]

这道计算同时检查几何末态、角动量比例和重力功的正号。物块在第3问向下走，而在这里小球向外、绳变短，物块向上走，两个问题的势能变化符号相反。若把第3问的
\(-mg\delta s\) 原封不动搬来，就会得到不正确的临界冲量。

\newpage

\section{第 2 题}\label{ux7b2c-2-ux9898}

某种物质的电导率是非线性的，其电阻率 \(\rho\) 与电场强度 \(E\) 满足

\[
\rho=\rho_0+AE^2,
\]

其中

\[
\rho_0=1.7\times10^7\,\Omega\cdot\mathrm m,
\qquad
A=1.0\times10^{-3}\,\Omega\cdot\mathrm m^3/\mathrm V^2.
\]

用该物质充满平行板电容器的两块极板之间，极板面积 \(S=1\,\mathrm{m^2}\)。

\begin{enumerate}
\def\labelenumi{\arabic{enumi}.}
\tightlist
\item
  求流过电容器的电流最大值 \(I_{\max}\)；
\item
  极板间距
  \(d=1\,\mathrm{cm}\)，在两极板间施加可变电压，求介质内最大焦耳热功率，并画出
  \(P-U\) 关系图；
\item
  保持电压 \(U_1=2.0\times10^3\,\mathrm V\) 不变，改变极板间距
  \(d\)。求最大功率、达到最大功率时的 \(d_1\)，并定性画出 \(P-d\) 图。
\end{enumerate}

\begin{figure}[H]
\centering
\includegraphics[width=0.9\linewidth,height=\textheight,keepaspectratio,alt={图 2-1　非线性导电介质填充的平行板电容器。}]{../figures/question-02.pdf}
\caption*{图 2-1　非线性导电介质填充的平行板电容器。}
\end{figure}

\subsection{一、题型定位}\label{ux4e00ux9898ux578bux5b9aux4f4d-1}

电阻率随电场增强而增大，所以电流与功率的极值并不在同一处。先写电流密度，再分别固定间距或电压，检查函数的有限峰值与渐近上限。

\subsection{二、知识点}\label{ux4e8cux77e5ux8bc6ux70b9-1}

\begin{enumerate}
\def\labelenumi{\arabic{enumi}.}
\tightlist
\item
  电流密度 \(j=\sigma E=E/\rho\)；
\item
  平行板中的近似匀强电场 \(E=U/d\)；
\item
  电流 \(I=jS\)；
\item
  焦耳功率 \(P=UI\)；
\item
  求导或均值不等式判断极值；
\item
  有限极值与渐近上限的区别。
\end{enumerate}

\subsection{三、物理图像}\label{ux4e09ux7269ux7406ux56feux50cf-1}

电场较弱时，\(\rho\approx\rho_0\)，故 \(j\) 随 \(E\)
近似线性增大；电场很强时，\(AE^2\) 主导，\(j\approx1/(AE)\)，反而随
\(E\) 减小。因此电流存在有限最大值。固定 \(d\) 时，功率 \(P\)
随电压增大后趋于常量；固定 \(U\) 时，\(d\)
过小和过大都会使功率趋近于零，因此存在有限最优间距。

\subsection{四、解题入口}\label{ux56dbux89e3ux9898ux5165ux53e3-1}

先写统一表达式

\[
j(E)=\frac{E}{\rho_0+AE^2}.
\]

第（1）问对 \(E\) 求极值。第（2）、（3）问把 \(E=U/d\) 代入后分别把
\(U\) 或 \(d\) 视为自变量。判断``最大值''时要看题目允许的自变量范围：若
\(U\) 可无限增大，\(P(U)\) 只有上确界而无有限电压处的极大点。

\subsection{五、解析}\label{ux4e94ux89e3ux6790-1}

\begin{figure}[H]
\centering
\includegraphics[width=0.9\linewidth,height=\textheight,keepaspectratio,alt={图 2-2　电流密度峰值与固定电压时的功率峰值。}]{../figures/solution-02.pdf}
\caption*{图 2-2　电流密度峰值与固定电压时的功率峰值。}
\end{figure}

\textbf{（1）最大电流}

电流密度为

\[
j(E)=\frac{E}{\rho_0+AE^2}.
\tag{1}
\]

求导：

\[
\begin{aligned}
\frac{\mathrm dj}{\mathrm dE}
&=\frac{(\rho_0+AE^2)-E(2AE)}
{(\rho_0+AE^2)^2}\\
&=\frac{\rho_0-AE^2}
{(\rho_0+AE^2)^2}.
\end{aligned}
\tag{2}
\]

令导数为零：

\[
\rho_0-AE^2=0,
\]

得到

\[
E_{\rm m}=\sqrt{\frac{\rho_0}{A}}.
\tag{3}
\]

当 \(E<E_{\rm m}\) 时导数为正，当 \(E>E_{\rm m}\)
时导数为负，因此该点确为最大值。代回式（1）：

\[
\begin{aligned}
j_{\max}
&=\frac{\sqrt{\rho_0/A}}
{\rho_0+A(\rho_0/A)}\\
&=\frac{\sqrt{\rho_0/A}}{2\rho_0}\\
&=\frac1{2\sqrt{\rho_0A}}.
\end{aligned}
\tag{4}
\]

所以

\[
I_{\max}=Sj_{\max}
=\frac{S}{2\sqrt{\rho_0A}}.
\tag{5}
\]

代入数据：

\[
\sqrt{\rho_0A}
=\sqrt{1.7\times10^7\times10^{-3}}
=\sqrt{1.7\times10^4}
\approx130.38.
\]

于是

\[
\boxed{I_{\max}\approx3.83\times10^{-3}\,\mathrm A
=3.83\,\mathrm{mA}}.
\tag{6}
\]

\textbf{（2）固定板间距时的 \(P-U\) 关系}

由 \(E=U/d\)，

\[
I
=S\frac{U/d}{\rho_0+A(U/d)^2}.
\]

分子、分母同时乘以 \(d\)：

\[
I=\frac{SU}{\rho_0d+AU^2/d}.
\tag{7}
\]

因此

\[
\begin{aligned}
P(U)
&=UI\\
&=\frac{SU^2}{\rho_0d+AU^2/d}\\
&=\frac{SdU^2}{\rho_0d^2+AU^2}.
\end{aligned}
\tag{8}
\]

把式（8）除以 \(U^2\) 的主导项来考察大电压极限：

\[
\lim_{U\to\infty}P(U)
=\frac{Sd}{A}.
\tag{9}
\]

并且

\[
\frac{\mathrm dP}{\mathrm dU}
=\frac{2Sd\rho_0d^2U}
{(\rho_0d^2+AU^2)^2}>0
\qquad(U>0).
\tag{10}
\]

所以 \(P\) 随 \(U\) 单调增加，只在 \(U\to\infty\)
时趋近上限。题目所谓最大焦耳热功率应理解为上确界，任何有限电压都达不到它：

\[
P_{\max}=\frac{Sd}{A}.
\]

代入 \(S=1\,\mathrm{m^2}\)、\(d=10^{-2}\,\mathrm m\)、\(A=10^{-3}\)：

\[
\boxed{P_{\max}=10\,\mathrm W}.
\tag{11}
\]

\(P-U\) 图从原点出发，初始近似二次增长，随后逐渐变平，以
\(P=10\,\mathrm W\) 为水平渐近线。

\textbf{（3）固定电压时的最优间距}

固定 \(U_1\)，由式（8）

\[
P(d)=\frac{SU_1^2d}{\rho_0d^2+AU_1^2}.
\tag{12}
\]

求导：

\[
\frac{\mathrm dP}{\mathrm dd}
=SU_1^2
\frac{AU_1^2-\rho_0d^2}
{(\rho_0d^2+AU_1^2)^2}.
\tag{13}
\]

令分子为零：

\[
AU_1^2-\rho_0d_1^2=0,
\]

所以

\[
\boxed{d_1=U_1\sqrt{\frac A{\rho_0}}}.
\tag{14}
\]

代入数据：

\[
\begin{aligned}
d_1
&=2.0\times10^3
\sqrt{\frac{1.0\times10^{-3}}{1.7\times10^7}}\\
&=2.0\times10^3\sqrt{5.882\times10^{-11}}\\
&\approx1.534\times10^{-2}\,\mathrm m.
\end{aligned}
\]

即

\[
\boxed{d_1\approx1.53\,\mathrm{cm}}.
\tag{15}
\]

在极值点有

\[
\rho_0d_1^2=AU_1^2.
\]

代入式（12）：

\[
\begin{aligned}
P_{\max}
&=\frac{SU_1^2d_1}{2AU_1^2}\\
&=\frac{Sd_1}{2A}\\
&=\frac{SU_1}{2\sqrt{\rho_0A}}.
\end{aligned}
\]

数值为

\[
\boxed{P_{\max}\approx7.67\,\mathrm W}.
\tag{16}
\]

当 \(d\to0\) 时 \(P\to0\)；当 \(d\to\infty\) 时
\(P\sim SU_1^2/(\rho_0d)\to0\)。因此 \(P-d\) 图只有一个峰值，峰的位置为
\(d_1\)。

\textbf{三条曲线使用同一场强尺度}

定义 \(E_*=\sqrt{\rho_0/A}\)。固定板距时取 \(x=U/(dE_*)=E/E_*\)，则

\[\frac I{I_{\max}}=\frac{2x}{1+x^2},\qquad \frac P{P_\infty}=\frac{x^2}{1+x^2},\qquad P_\infty=Sd/A.\]

固定电压而改变板距，取 \(y=d/d_*\)，\(d_*=U\sqrt{A/\rho_0}\)，则

\[\frac P{P_{\max}}=\frac{2y}{1+y^2}.\]

解析图分别画出这三条曲线。第一、三条在无量纲坐标1处达到峰值；第二条在1处只有上限的一半，随后持续增加。三者都由同一材料关系产生，但``电压增大''和``板距减小''同时改变电场及几何因子，不能互相替代。

在低场 \(AE^2\ll\rho_0\)，\(I\propto U\)、\(P\propto U^2\)；在高场
\(AE^2\gg\rho_0\)，\(I\propto U^{-1}\)、\(P\to P_\infty\)。这两个极限还可独立核对图像两端。

\textbf{负微分电导并不等于负电阻}

材料给 \(j(E)=E/(\rho_0+AE^2)\)，故

\[\frac{dj}{dE}=\frac{\rho_0-AE^2}{(\rho_0+AE^2)^2}.\]

当 \(E>E_*\)，电流随电压增加而减小，微分电导为负。但静态比值
\(E/j=\rho_0+AE^2\) 始终为正，焦耳功率 \(jE=E^2/(\rho_0+AE^2)\)
也始终为正。两种``电阻''的定义不同，这可解释图中电流下降而功率上升的现象。

若电压由理想电源固定，工作点由给定 \(U\)
唯一决定；若串接外电阻或采用电流源，工作点的稳定性还与外电路相关，不能仅凭这条材料曲线判断。题目只求静态关系，不包含温度随发热改变材料参数的反馈。要研究实际持续工作，应补充散热条件，不能把渐近功率理解成装置一定能安全维持的额定功率。

计算最优板距时还有一个短检查：\(d_*=U/E_*\)，因此改变电压会同比改变最优距离，但最优材料电场不变。\(S\)
只把电流和功率同比放大，不改变曲线峰值的位置。

\textbf{从图上反求工作区间}

用 \(E_* =\sqrt{\rho_0/A}\) 定义 \(x=E/E_*\)，图中第一条曲线满足

\[\frac I{I_{\max}}=\frac{2x}{1+x^2}.\]

若要求电流至少达到峰值的 \(80\%\)，不等式为
\(2x/(1+x^2)\ge0.8\)。移项并因式分解：

\[4x^2-10x+4\le0,\qquad (2x-1)(x-2)\le0.\]

所以
\(1/2\le x\le2\)。同一个低于峰值的电流可以对应两个电场：一个在增长支，一个在下降支。两根乘积为1，正好对应曲线在
\(x\leftrightarrow1/x\)
下不变的性质。若只给电流读数，不给电压或电场，就不能辨认介质处在哪一支。

题给参数对应 \(E_*=1.30384\times10^5\,\mathrm{V/m}\)。当间距
\(d=1.00\,\mathrm{cm}\) 时，峰值电流电压为

\[U_*=dE_*=1.30384\times10^3\,\mathrm V.\]

电流达到 \(80\%\) 峰值的电压区间因此约为 \(652\,\mathrm V\) 到
\(2.608\,\mathrm{kV}\)。但这一电流峰值对应的焦耳功率只有固定间距功率上限的一半：

\[P(U_*)=U_*I_{\max}=\frac{Sd}{2A}=5.00\,\mathrm W.\]

这可在图中把第一条曲线的最高点与第二条曲线的半高点对应起来。电流越过峰值后开始下降，功率仍可继续增大，因为外加电压增加得更快。

固定电压 \(U_1\) 时，第三条曲线用 \(y=d/d_*\) 也有同一个函数
\(2y/(1+y^2)\)。功率达到最优值的 \(80\%\)，要求
\(d_*/2\le d\le2d_*\)。代入 \(d_*=1.53393\,\mathrm{cm}\)，得到约
\(0.767\) 到
\(3.068\,\mathrm{cm}\)。这两个区间分别控制不同自变量，不能把横轴数值混用。若设计还受击穿场强
\(E\le E_b\) 限制，应将允许区间与 \(d\ge U_1/E_b\)
再取交集；只有当最优点仍在允许区间内，正文的无约束峰值才能实际取到。

\textbf{加入场强限制后的具体最优点}

取前面的固定电压 \(U_1=2000\,\mathrm V\)，再明确增加场强限制
\(E_b=1.00\times10^5\,\mathrm{V/m}\)。允许间距必须满足

\[d\ge d_{\min}=U_1/E_b=2.00\,\mathrm{cm}.\]

无约束峰值在 \(d_*=1.53393\,\mathrm{cm}\)，落在允许区间之外。因
\(d>d_*\) 后功率随间距单调下降，受约束最大值只能在左端点取得：

\[P_{\rm allowed}=\frac{SU_1^2d_{\min}}{\rho_0d_{\min}^2+AU_1^2}=\frac{80000}{10800}\,\mathrm W\simeq7.407\,\mathrm W.\]

对应电流 \(I=P/U_1\simeq3.704\,\mathrm{mA}\)，小于材料本身的峰值电流
\(3.835\,\mathrm{mA}\)。把这个点放回第三条图线，横坐标为
\(d_{\min}/d_*=1.30384\)，位于峰的右侧，纵坐标为约
\(0.9658\)。图上的位置、求导的符号和直接代入数值三者相符。

这只是增加场强参数后的算例，原题没有给击穿场强，原题答案仍为无约束结果。处理带额外限制的设计问题时，顺序应是先求候选驻点，再与允许边界比较，最后将选中的点代回原函数。不能看到驻点就立刻宣称它是可达到的最大功率。

\newpage

\section{第 3 题}\label{ux7b2c-3-ux9898}

有一根质量分布不均匀的绳子，总长为 \(L\)。绳上各点的线密度
\(\lambda(x)\) 与该点到左端点 \(P\)
的距离成正比。把绳放在水平桌面上，使长度为 \(l_0\)
的一段从光滑桌角竖直下垂，桌高大于绳长。由静止释放。绳与桌面间的动摩擦因数为
\(\mu\)，桌角光滑，不考虑绳在运动中飞起。

\begin{enumerate}
\def\labelenumi{\arabic{enumi}.}
\tightlist
\item
  已知 \(\lambda(x)=kx\)。若要使绳释放后不动，至少需要多大的静摩擦因数
  \(\mu'\)；
\item
  若随后摩擦因数降为
  \(\mu<\mu'\)，从初始状态到绳刚离开桌面，摩擦力对绳做的总功是多少；
\item
  从初始状态到绳刚离开桌面，绳的重力势能减少量是多少；
\item
  绳刚离开桌面时的速度是多少。
\end{enumerate}

\begin{figure}[H]
\centering
\includegraphics[width=0.717\linewidth,height=\textheight,keepaspectratio,alt={图 3-1　线密度随位置增大的绳子跨过光滑桌角。}]{../figures/question-03.pdf}
\caption*{图 3-1　线密度随位置增大的绳子跨过光滑桌角。}
\end{figure}

\subsection{一、题型定位}\label{ux4e00ux9898ux578bux5b9aux4f4d-2}

绳不断离开桌面，悬垂部分的重量和剩余摩擦力同时改变。用材料坐标跟踪每个绳元，积分重力功和摩擦功，再判断能否从静止开始滑动。

\subsection{二、知识点}\label{ux4e8cux77e5ux8bc6ux70b9-2}

\begin{enumerate}
\def\labelenumi{\arabic{enumi}.}
\tightlist
\item
  线密度积分求质量；
\item
  最大静摩擦与动摩擦；
\item
  摩擦力随位移变化时的积分功；
\item
  连续质量分布的重力势能；
\item
  整根绳速度相同的不可伸长约束；
\item
  功能关系求末速度。
\end{enumerate}

\subsection{三、物理图像}\label{ux4e09ux7269ux7406ux56feux50cf-2}

坐标 \(x\) 从左端点 \(P\) 沿绳向右量。初态在桌面上的绳段为
\(0\le x\le L-l_0\)，下垂段为 \(L-l_0\le x\le L\)。由于
\(\lambda=kx\)，靠近自由端的下垂绳更重，这正是绳具有下滑趋势的原因。

当绳沿桌面移动距离 \(s\) 时，仍留在桌面上的原绳段为
\(0\le x\le L-l_0-s\)。其质量随 \(s\) 递减，故动摩擦力不是常量。

\subsection{四、解题入口}\label{ux56dbux89e3ux9898ux5165ux53e3-2}

第（1）问比较初态下垂段重力与桌面段最大静摩擦。第（2）问以绳的位移 \(s\)
为积分变量，先求任意 \(s\)
时桌面段的支持力。第（3）问分别计算初态和末态势能，统一把桌面高度取为零。第（4）问用

\[
K_f+\Delta E_p=W_f
\]

或``势能减少量加摩擦功等于末动能''求速度。

\subsection{五、解析}\label{ux4e94ux89e3ux6790-2}

\begin{figure}[H]
\centering
\includegraphics[width=0.837\linewidth,height=\textheight,keepaspectratio,alt={图 3-2　位移过程中的桌面绳段与刚离桌时的竖直绳。}]{../figures/solution-03.pdf}
\caption*{图 3-2　位移过程中的桌面绳段与刚离桌时的竖直绳。}
\end{figure}

\textbf{（1）最小静摩擦因数}

初态桌面段质量为

\[
\begin{aligned}
m_{\rm t}
&=\int_0^{L-l_0}\lambda(x)\,\mathrm dx\\
&=\int_0^{L-l_0}kx\,\mathrm dx\\
&=\frac12k(L-l_0)^2.
\end{aligned}
\tag{1}
\]

桌面对绳的总支持力为

\[
N=m_{\rm t}g
=\frac12kg(L-l_0)^2.
\tag{2}
\]

下垂段质量为

\[
\begin{aligned}
m_{\rm h}
&=\int_{L-l_0}^{L}kx\,\mathrm dx\\
&=\frac12k\left[L^2-(L-l_0)^2\right]\\
&=\frac12k(2L-l_0)l_0.
\end{aligned}
\tag{3}
\]

下垂段重力为

\[
G_{\rm h}
=\frac12kg(2L-l_0)l_0.
\tag{4}
\]

刚好不滑时，最大静摩擦等于下垂段重力：

\[
\mu'N=G_{\rm h}.
\]

代入式（2）、（4）并约去 \(kg/2\)：

\[
\mu'(L-l_0)^2=(2L-l_0)l_0.
\]

所以

\[
\boxed{\mu'_{\min}
=\frac{(2L-l_0)l_0}{(L-l_0)^2}}.
\tag{5}
\]

\textbf{（2）摩擦力的总功}

设绳已经移动距离 \(s\)。此时仍在桌面上的绳对应
\(0\le x\le L-l_0-s\)，其质量为

\[
m_{\rm t}(s)
=\int_0^{L-l_0-s}kx\,\mathrm dx
=\frac12k(L-l_0-s)^2.
\tag{6}
\]

动摩擦力大小为

\[
f(s)=\mu m_{\rm t}(s)g
=\frac12\mu kg(L-l_0-s)^2.
\tag{7}
\]

摩擦力与位移方向相反，所以微元功

\[
\mathrm dW_f=-f(s)\,\mathrm ds.
\]

绳刚离开桌面时，移动总距离为 \(L-l_0\)。因此

\[
\begin{aligned}
W_f
&=-\frac12\mu kg
\int_0^{L-l_0}(L-l_0-s)^2\,\mathrm ds.
\end{aligned}
\]

令 \(u=L-l_0-s\)，则 \(\mathrm du=-\mathrm ds\)。当 \(s=0\) 时
\(u=L-l_0\)；当 \(s=L-l_0\) 时 \(u=0\)。于是

\[
\begin{aligned}
W_f
&=-\frac12\mu kg
\int_{L-l_0}^{0}u^2(-\mathrm du)\\
&=-\frac12\mu kg
\int_0^{L-l_0}u^2\,\mathrm du\\
&=-\frac12\mu kg
\left[\frac{u^3}{3}\right]_0^{L-l_0}\\
&=-\frac16\mu kg(L-l_0)^3.
\end{aligned}
\]

所以

\[
\boxed{W_f=-\frac16\mu kg(L-l_0)^3}.
\tag{8}
\]

\textbf{（3）重力势能减少量}

取桌面高度为零，向下为正深度。初态中，位置标号为 \(x\)
的下垂微元位于桌面下

\[
z=x-(L-l_0).
\]

微元质量为 \(\mathrm dm=kx\,\mathrm dx\)，其势能为

\[
\mathrm dE_{p1}
=-gz\,\mathrm dm
=-kg[x-(L-l_0)]x\,\mathrm dx.
\]

积分得初态势能：

\[
\begin{aligned}
E_{p1}
&=-kg\int_{L-l_0}^{L}
[x-(L-l_0)]x\,\mathrm dx\\
&=-kg\int_{L-l_0}^{L}
\left[x^2-(L-l_0)x\right]\mathrm dx.
\end{aligned}
\]

逐项积分：

\[
\begin{aligned}
E_{p1}
&=-kg\left[
\frac{x^3}{3}-\frac{L-l_0}{2}x^2
\right]_{L-l_0}^{L}\\
&=-kg\left(\frac12Ll_0^2-\frac16l_0^3\right).
\end{aligned}
\tag{9}
\]

刚离开桌面时整根绳竖直，标号 \(x\) 的微元位于桌面下深度
\(x\)。末态势能为

\[
\begin{aligned}
E_{p2}
&=-kg\int_0^Lx^2\,\mathrm dx\\
&=-\frac13kgL^3.
\end{aligned}
\tag{10}
\]

势能减少量为 \(E_{p1}-E_{p2}\)：

\[
\begin{aligned}
\Delta U_{\rm drop}
&=kg\left(
\frac13L^3-\frac12Ll_0^2+\frac16l_0^3
\right)\\
&=\frac16kg
\left(2L^3-3Ll_0^2+l_0^3\right).
\end{aligned}
\]

因此

\[
\boxed{\Delta U_{\rm drop}
=\frac16kg
\left(2L^3-3Ll_0^2+l_0^3\right)}.
\tag{11}
\]

\textbf{（4）刚离桌时的速度}

整根绳的总质量为

\[
m_{\rm tot}
=\int_0^Lkx\,\mathrm dx
=\frac12kL^2.
\tag{12}
\]

绳不可伸长，各微元沿绳方向的速率相同，故末动能

\[
K_f
=\frac12m_{\rm tot}v^2
=\frac14kL^2v^2.
\tag{13}
\]

初态静止。势能减少量中一部分转化为动能，另一部分被摩擦耗散：

\[
K_f=\Delta U_{\rm drop}+W_f.
\tag{14}
\]

代入式（8）、（11）、（13）：

\[
\frac14kL^2v^2
=\frac16kg
\left[
2L^3-3Ll_0^2+l_0^3-\mu(L-l_0)^3
\right].
\]

约去 \(k\)，两边乘以 \(4/L^2\)：

\[
v^2
=\frac{2g}{3L^2}
\left[
2L^3-3Ll_0^2+l_0^3-\mu(L-l_0)^3
\right].
\]

所以

\[
\boxed{
v=\frac1L
\sqrt{\frac{2g}{3}
\left[
2L^3-3Ll_0^2+l_0^3-\mu(L-l_0)^3
\right]}
}.
\tag{15}
\]

\textbf{启动后为什么不会停下}

下落位移为 \(s\) 时，桌面剩余绳长 \(L-l_0-s\)。沿整根绳的净驱动力为

\[F(s)=\frac{kg}{2}\left[L^2-(1+\mu)(L-l_0-s)^2\right].\]

对 \(0\le s\le L-l_0\)，\(F'(s)=kg(1+\mu)(L-l_0-s)\ge0\)。只要初态
\(F(0)>0\)，以后就一直为正。又因为整根绳的质量固定为 \(kL^2/2\)，有

\[\frac12M v^2(s)=\int_0^sF(\xi)\,d\xi.\]

这既证明全过程速度增长，也在终点还原正文能量结果。由终点动能非负所得摩擦上界比较宽，不能代替初态启动上界。若
\(l_0=0\) 且完全静止、没有扰动，初态不具备悬垂驱动力；式中 \(l_0\to0^+\)
的极限不能解释为这种精确初态会自行运动。

\textbf{任意下落位置的速度，而不只终点}

整绳质量 \(M=kL^2/2\)。把前面的净力从0积分到 \(s\)，有

\[v^2(s)=2gs-\frac{2g(1+\mu)}{3L^2}\left[(L-l_0)^3-(L-l_0-s)^3\right].\]

微分检查 \(\tfrac12M\,d(v^2)/ds=F(s)\)，可避免摩擦功积分的上下限错误。取
\(s=L-l_0\) 后，展开前一项
\(3L^2(L-l_0)\)，便能与正文终点结果逐项比较。这里所有绳元速度大小相同，绳在桌边改变方向；不能把桌面部分视为不动而只给悬垂部分写动能。

若 \(\mu=\mu'\)
恰在初态临界，\(F(0)=0\)；一旦给很小的正位移，\(F(s)>0\)，所以这是对扰动不稳定的平衡。是否会由完全静止自行启动，取决于静摩擦临界与初始扰动，不能用终点正动能回答这个初态问题。原题将静、动摩擦取同一数值时，严格启动条件写成
\(\mu<\mu'\) 最清楚。

材料坐标 \(x\) 表示沿原绳从较轻端到较重端的位置，状态位移 \(s\)
表示已经下落多远。两者的单位相同，但积分对象不同；图中将原绳坐标和悬垂长度分别标注，才能核对
\(\lambda=kx\) 的质量元。

\textbf{直接积分每个绳元的重力势能}

初始悬垂段对应材料坐标 \(x\in[L-l_0,L]\)，其高度相对桌面为
\(y=-(x-L+l_0)\)。终态整绳竖直，材料坐标 \(x\) 的高度为 \(-x\)。于是

\[U_i=-kg\int_{L-l_0}^{L}x(x-L+l_0)dx,\qquad U_f=-kg\int_0^Lx^2dx.\]

前一个积分取 \(u=x-L+l_0\)，可写为

\[U_i=-kg\left[\frac{(L-l_0)l_0^2}{2}+\frac{l_0^3}{3}\right],\qquad U_f=-\frac{kgL^3}{3}.\]

所以重力释放能量为

\[U_i-U_f=\frac{kg}{6}(2L^3-3Ll_0^2+l_0^3).\]

它恰好等于正文重力势能差。这个推导从每个材料微元的初末高度开始，不依赖悬垂段质心的位置，可作为另一条检查路线。摩擦功仍只来自桌面部分，不能再为已经离开桌面的微元补一段摩擦。

取 \(l_0=L\)，初始本来已全部竖直，两积分相同，释放能量为零；取
\(l_0\to0^+\)，初始势能趋零，释放能量趋
\(kgL^3/3\)。两个端点分别检查符号和三次长度因子。

\newpage

\section{第 4 题}\label{ux7b2c-4-ux9898}

两个相同的线轴对称放置在一个对称山坡的两侧，山坡顶角为
\(150^\circ\)。线轴内筒绕有不可伸长的轻绳，轻绳与内筒之间不打滑，内筒半径是外筒半径的一半。绳子中点
\(P\)
从静止开始竖直向上移动。山坡的摩擦因数足够大，线轴始终做无滑动纯滚动。求当
\(P\) 点速度为 \(v\)、两段绳的夹角 \(2\alpha=60^\circ\)
时，两线轴轴线沿水平方向的接近速度 \(u\)。

\begin{figure}[H]
\centering
\includegraphics[width=0.515\linewidth,height=\textheight,keepaspectratio,alt={图 4-1　对称山坡、线轴与上移绳中点。}]{../figures/question-04.pdf}
\caption*{图 4-1　对称山坡、线轴与上移绳中点。}
\end{figure}

\subsection{一、题型定位}\label{ux4e00ux9898ux578bux5b9aux4f4d-3}

两只线轴同时受坡面纯滚动和内轴不滑绳约束。把切点速度投影到绳方向，再把轴心速度投影到水平方向，两次投影的角度不同。

\subsection{二、知识点}\label{ux4e8cux77e5ux8bc6ux70b9-3}

\begin{enumerate}
\def\labelenumi{\arabic{enumi}.}
\tightlist
\item
  纯滚动条件 \(v_C=\omega R\)；
\item
  刚体上一点速度 \(\bm v_B=\bm v_C+\bm\omega\times\bm r_{CB}\)；
\item
  不可伸长绳两端沿绳方向的速度分量关系；
\item
  对称运动与水平接近速度；
\item
  山坡顶角、坡角和绳角之间的几何换算。
\end{enumerate}

\subsection{三、物理图像}\label{ux4e09ux7269ux7406ux56feux50cf-3}

山坡顶角为 \(150^\circ\)，因此每侧坡面与水平面的夹角为

\[
\delta=\frac{180^\circ-150^\circ}{2}=15^\circ.
\]

当 \(2\alpha=60^\circ\) 时，每段绳与竖直方向夹角为
\(\alpha=30^\circ\)，也就是与水平夹角
\(60^\circ\)。左侧线轴沿坡面向右上方靠近山顶，其轴心速度与水平夹角为
\(15^\circ\)。轴心速度方向与绳方向夹角因此为

\[
\beta=60^\circ-15^\circ=45^\circ.
\]

\subsection{四、解题入口}\label{ux56dbux89e3ux9898ux5165ux53e3-3}

只研究左侧线轴。设外筒半径为 \(2r\)、内筒半径为 \(r\)，轴心沿坡速度为
\(v_0\)。先由纯滚动求角速度 \(\omega\)，再把绕线点速度投影到绳方向，与
\(P\) 点速度在绳方向的投影相等。最后把一个线轴的水平速度乘以
\(2\)，得到两轴接近速度。

\subsection{五、解析}\label{ux4e94ux89e3ux6790-3}

\begin{figure}[H]
\centering
\includegraphics[width=0.45\linewidth,height=\textheight,keepaspectratio,alt={图 4-2　外筒接触点、内筒绕线点与速度合成。}]{../figures/solution-04.pdf}
\caption*{图 4-2　外筒接触点、内筒绕线点与速度合成。}
\end{figure}

设左线轴外半径为 \(R=2r\)，内半径为 \(r\)。轴心沿坡向山顶的速度为
\(v_0\)。

\textbf{（1）纯滚动条件}

线轴与坡面接触点瞬时速度为零，因此

\[
v_0=\omega R=2r\omega.
\tag{1}
\]

所以

\[
\omega r=\frac{v_0}{2}.
\tag{2}
\]

\textbf{（2）不可伸长绳的速度投影}

取内筒上绳的切点为 \(B\)。绳与内筒不打滑，所以绳在该点的速度等于 \(B\)
点沿绳方向的速度。

\(P\) 点速度竖直向上，绳与竖直夹角为 \(\alpha\)，故 \(P\)
点速度沿绳方向的分量为

\[
v\cos\alpha.
\tag{3}
\]

轴心速度与绳方向夹角为 \(\beta\)，其沿绳分量为

\[
v_0\cos\beta.
\tag{4}
\]

内筒转动在切点产生沿绳方向的速度
\(\omega r\)。在图示卷绕方向下，它与式（4）同向。因此

\[
v\cos\alpha=v_0\cos\beta+\omega r.
\tag{5}
\]

代入式（2）：

\[
v\cos\alpha
=v_0\left(\cos\beta+\frac12\right).
\]

解得

\[
v_0
=\frac{v\cos\alpha}{\cos\beta+1/2}.
\tag{6}
\]

当 \(\alpha=30^\circ\)、\(\beta=45^\circ\) 时，

\[
v_0
=\frac{v\cos30^\circ}{\cos45^\circ+1/2}.
\tag{7}
\]

\textbf{（3）两轴的水平接近速度}

单个线轴轴心的水平分速度为

\[
v_{0x}=v_0\cos15^\circ.
\]

左右两线轴对称靠近，因此轴线之间距离的减小速率为

\[
u=2v_0\cos15^\circ.
\tag{8}
\]

把式（7）代入：

\[
u
=\frac{2v\cos30^\circ\cos15^\circ}
{\cos45^\circ+1/2}.
\tag{9}
\]

利用

\[
\cos15^\circ=\frac{\sqrt6+\sqrt2}{4},
\quad
\cos30^\circ=\frac{\sqrt3}{2},
\quad
\cos45^\circ=\frac{\sqrt2}{2},
\]

分子为

\[
2\cos30^\circ\cos15^\circ
=\frac{\sqrt3(\sqrt6+\sqrt2)}4
=\frac{3\sqrt2+\sqrt6}{4}.
\]

分母为

\[
\cos45^\circ+\frac12
=\frac{\sqrt2+1}{2}.
\]

整理可得

\[
\boxed{
u=\frac{3+\sqrt3}{2+\sqrt2}\,v
}.
\tag{10}
\]

数值上

\[
\boxed{u\approx1.386\,v}.
\]

\textbf{绕线点的速度要投影}

线轴沿 \(15^\circ\)
斜面向右上滚动，轴心速度沿坡面，角速度沿顺时针方向。内轴半径为
\(r\)，外轮半径为 \(2r\)，纯滚动给 \(\omega=v_C/(2r)\)。绳在绕线点沿
\(60^\circ\) 方向离开，因而轴心平动的绳向分量是
\(v_C\cos45^\circ\)，绕心转动的分量是 \(\omega r\)，两者相加。

但绕线点的完整速度未必沿绳方向。不可伸长条件只要求它与绳接触点在绳方向上的速度相等，不能把两个不共线的速度大小直接相加。解析图将轴心速度、绕心速度及绳方向投影分别标明，再与顶端
\(P\) 点速度的 \(\cos30^\circ\) 投影比较。

用瞬时中心 \(A\)
核对时，\(\bm v_B=\bm\omega\times\overrightarrow{AB}\)；先取点积
\(\bm v_B\cdot\hat{\bm t}_{\rm rope}\)，再与绳速比较，得到同一关系。

\textbf{换内轴半径后的通用速度关系}

保留同一几何，令内、外轴半径分别为 \(r,R\)。绕心速度在绳向的分量为
\(\omega r\)，纯滚动给 \(\omega=v_C/R\)。顶端点 \(P\)
的速度在绳向的投影为 \(v_P\cos30^\circ\)，所以

\[v_P\cos30^\circ=v_C\left(\cos45^\circ+\frac rR\right).\]

由装置左右对称，轴心的两侧运动对 \(P\)
的绳长变化均有贡献；正文定义的拉绳速度 \(u=2v_C\cos15^\circ\)，因此

\[\frac uv=\frac{2\cos30^\circ\cos15^\circ}{\cos45^\circ+r/R}.\]

取 \(r/R=1/2\) 可还原
\((3+\sqrt3)/(2+\sqrt2)\)。若内轴半径趋零，绕心作用消失，速度仍由轴心绳向投影决定；若改变绳的切点到另一侧，转动项符号会变，不能只将半径换数值。题图的切点位置正是决定正负号的信息。

这一关系只涉及运动学，不能据此求需要多大的拉力。若题目继续问加速度或摩擦，必须再给质量、转动惯量和外力，并检查静摩擦是否足以维持纯滚动；``不打滑''在本问作为已给约束使用。

\textbf{把每个速度投影回代到同一瞬间}

令绳中点上升速度 \(v=1.00\,\mathrm{m/s}\)，由正文得左线轴轴心沿坡速度

\[v_C=\frac{\cos30^\circ}{\cos45^\circ+1/2}=0.71744\,\mathrm{m/s}.\]

它沿绳的平动分量为 \(0.50731\,\mathrm{m/s}\)，内轴转动贡献为
\(\omega r=v_C/2=0.35872\,\mathrm{m/s}\)。两者相加等于
\(0.86603\,\mathrm{m/s}\)，正是绳中点速度在该段绳方向的投影：

\[0.50731+0.35872=\cos30^\circ.\]

轴心的水平速度则为
\(v_C\cos15^\circ=0.69301\,\mathrm{m/s}\)，两轴接近速度是其两倍，即
\(1.3860\,\mathrm{m/s}\)。这个回代同时检查了两个不同投影角：\(45^\circ\)
用于轴心速度投影到绳方向，\(15^\circ\)
用于投影到水平方向。二者没有互换的余地。

图中切点的内轴半径垂直于绳；外轮接触半径垂直于坡面。取左轴顺时针角速度，内轴切点的转动速度沿绳向上，而外轮接触点的转动速度恰与轴心速度相反。先用这两个垂直关系检查图，再列速度约束，就能避免把不相切的示意线代入正确公式。

\newpage

\section{第 5 题}\label{ux7b2c-5-ux9898}

在自由空间中，半径为 \(R_0\) 的圆内接正方形四个顶点上各有一个质量为
\(m\) 的质点。其中两个带电 \(+q\)，另两个带电
\(-q\)，相邻顶点电荷异号。初始时四个质点以相同速率沿圆的切线顺时针运动。已知运动过程中，任一质点到圆心
\(O\) 的最小距离为 \(R_1<R_0\)；任意时刻四个电荷仍位于以 \(O\)
为中心的正方形四个顶点。忽略重力，只考虑静电作用。

\begin{enumerate}
\def\labelenumi{\arabic{enumi}.}
\tightlist
\item
  每个粒子的运动轨迹是什么曲线；
\item
  求一个粒子从初始位置运动到距圆心 \(R_1\) 处所需的时间；
\item
  在初态沿两个负电荷所在直线方向、距 \(O\) 点 \(l\gg R_0\)
  处，求电场强度精确到首个非零小量。
\end{enumerate}

\begin{figure}[H]
\centering
\includegraphics[width=0.608\linewidth,height=\textheight,keepaspectratio,alt={图 5-1　相邻异号的四电荷正方形与初始切向速度。}]{../figures/question-05.pdf}
\caption*{图 5-1　相邻异号的四电荷正方形与初始切向速度。}
\end{figure}

\subsection{一、题型定位}\label{ux4e00ux9898ux578bux5b9aux4f4d-4}

交替电荷的正方形可以保持形状，但各顶点的合力要先算清。把整体呼吸运动化成单个顶点的中心力问题，用能量和角动量求两次转向半径。

\subsection{二、知识点}\label{ux4e8cux77e5ux8bc6ux70b9-4}

\begin{enumerate}
\def\labelenumi{\arabic{enumi}.}
\tightlist
\item
  正方形构型中的库仑力矢量合成；
\item
  反平方中心力的开普勒椭圆轨道；
\item
  椭圆近心点、远心点与半长轴；
\item
  开普勒第三定律；
\item
  二项式展开到二阶；
\item
  电场远场展开与对称性消项。
\end{enumerate}

\subsection{三、物理图像}\label{ux4e09ux7269ux7406ux56feux50cf-4}

取上方的 \(+q\) 为研究对象。它受左右两个 \(-q\) 的吸引，合力沿半径指向
\(O\)；受下方 \(+q\) 的排斥，方向沿半径背离 \(O\)。合力仍指向
\(O\)，且大小与 \(1/R^2\) 成正比。于是每个粒子都等效为在固定中心 \(O\)
的反平方吸引力场中运动。

初始位置 \(R_0\) 是远心距，最短距离 \(R_1\)
是近心距。四条轨道只是彼此转过 \(90^\circ\)，每条都是以 \(O\)
为一个焦点的椭圆。

\subsection{四、解题入口}\label{ux56dbux89e3ux9898ux5165ux53e3-4}

先求单个粒子的径向合力，得到等效``引力参数''。第（1）问据此识别开普勒椭圆。第（2）问由远心距与近心距求半长轴
\(a=(R_0+R_1)/2\)，从远心点到近心点恰为半个周期。第（3）问写出轴线上四个电荷电场的精确表达式，再令
\(\varepsilon=R_0/l\) 做二阶展开。

\subsection{五、解析}\label{ux4e94ux89e3ux6790-4}

\begin{figure}[H]
\centering
\includegraphics[width=0.9\linewidth,height=\textheight,keepaspectratio,alt={图 5-2　单个电荷的径向合力与远场展开方向。}]{../figures/solution-05.pdf}
\caption*{图 5-2　单个电荷的径向合力与远场展开方向。}
\end{figure}

记库仑常量

\[
k_e=\frac1{4\pi\varepsilon_0}.
\]

\textbf{（1）粒子轨迹}

任一时刻正方形外接圆半径记为 \(R\)。相邻两顶点距离为

\[
\sqrt2R,
\]

对角两顶点距离为

\[
2R.
\]

研究上方 \(+q\)。它受一个相邻 \(-q\) 的吸引力大小为

\[
F_1=\frac{k_eq^2}{(\sqrt2R)^2}
=\frac{k_eq^2}{2R^2}.
\]

该力与指向 \(O\) 的半径夹角为
\(45^\circ\)，两个相邻负电荷的切向分量互相抵消，径向分量相加：

\[
\begin{aligned}
F_{\rm att}
&=2F_1\cos45^\circ\\
&=2\cdot\frac{k_eq^2}{2R^2}\cdot\frac1{\sqrt2}\\
&=\frac{k_eq^2}{\sqrt2R^2}.
\end{aligned}
\tag{1}
\]

对角位置的 \(+q\) 产生排斥力，方向背离 \(O\)，大小为

\[
F_{\rm rep}
=\frac{k_eq^2}{(2R)^2}
=\frac{k_eq^2}{4R^2}.
\tag{2}
\]

净力指向 \(O\)：

\[
\begin{aligned}
F_R
&=F_{\rm att}-F_{\rm rep}\\
&=\frac{k_eq^2}{R^2}
\left(\frac1{\sqrt2}-\frac14\right)\\
&=\frac{k_eq^2}{R^2}\frac{2\sqrt2-1}{4}.
\end{aligned}
\tag{3}
\]

定义

\[
\Gamma=\frac{k_eq^2(2\sqrt2-1)}4.
\tag{4}
\]

则运动方程为

\[
m\ddot{\bm r}
=-\frac{\Gamma}{R^3}\bm r.
\tag{5}
\]

这是标准反平方中心力。因初速度沿切线且
\(R_1>0\)，角动量不为零，轨道为椭圆，\(O\) 位于椭圆的一个焦点。因此

\[
\boxed{\text{每个粒子沿以 }O\text{ 为一个焦点的椭圆运动}}.
\tag{6}
\]

四条椭圆轨道形状相同，方位相差 \(90^\circ\)。

\textbf{（2）从 \(R_0\) 到 \(R_1\) 的时间}

初始半径最大，为椭圆远心距：

\[
r_{\rm a}=R_0.
\]

最短半径为近心距：

\[
r_{\rm p}=R_1.
\]

椭圆半长轴满足

\[
r_{\rm a}+r_{\rm p}=2a,
\]

故

\[
a=\frac{R_0+R_1}{2}.
\tag{7}
\]

把式（5）与万有引力形式

\[
\frac{GMm}{R^2}
\]

比较，可将等效的 \(GMm\) 替换为 \(\Gamma\)。开普勒第三定律给出周期

\[
T=2\pi\sqrt{\frac{ma^3}{\Gamma}}.
\tag{8}
\]

从远心点到近心点走过半个轨道周期，因此

\[
t=\frac T2
=\pi\sqrt{\frac{ma^3}{\Gamma}}.
\tag{9}
\]

代入式（4）、（7）：

\[
\begin{aligned}
t
&=\pi
\sqrt{
\frac{m[(R_0+R_1)/2]^3}
{k_eq^2(2\sqrt2-1)/4}
}\\
&=\pi
\sqrt{
\frac{m(R_0+R_1)^3}
{2k_eq^2(2\sqrt2-1)}
}.
\end{aligned}
\]

所以

\[
\boxed{
t=\pi\sqrt{
\frac{m(R_0+R_1)^3}
{2k_eq^2(2\sqrt2-1)}
}
}.
\tag{10}
\]

\textbf{（3）负电荷方向远处的电场}

取 \(x\) 轴沿两个负电荷方向，场点 \(P\) 位于 \(x=l\)，且
\(l\gg R_0\)。规定 \(+x\) 方向为正。

右侧 \(-q\) 位于 \(x=R_0\)，其电场指向左：

\[
E_{-R}=-\frac{k_eq}{(l-R_0)^2}.
\]

左侧 \(-q\) 位于 \(x=-R_0\)，其电场也指向左：

\[
E_{-L}=-\frac{k_eq}{(l+R_0)^2}.
\]

上、下两个 \(+q\) 的竖直分量互相抵消，水平分量相加。单个水平分量为

\[
\frac{k_eq}{l^2+R_0^2}
\frac{l}{\sqrt{l^2+R_0^2}},
\]

因此二者合计

\[
E_+=\frac{2k_eql}{(l^2+R_0^2)^{3/2}}.
\]

总场为

\[
E_x
=-\frac{k_eq}{(l-R_0)^2}
-\frac{k_eq}{(l+R_0)^2}
+\frac{2k_eql}{(l^2+R_0^2)^{3/2}}.
\tag{11}
\]

令

\[
\varepsilon=\frac{R_0}{l}\ll1.
\]

把式（11）提出 \(k_eq/l^2\)：

\[
E_x=\frac{k_eq}{l^2}
\left[
-(1-\varepsilon)^{-2}
-(1+\varepsilon)^{-2}
+2(1+\varepsilon^2)^{-3/2}
\right].
\tag{12}
\]

逐项展开到二阶：

\[
(1-\varepsilon)^{-2}
=1+2\varepsilon+3\varepsilon^2+O(\varepsilon^3),
\]

\[
(1+\varepsilon)^{-2}
=1-2\varepsilon+3\varepsilon^2+O(\varepsilon^3),
\]

\[
(1+\varepsilon^2)^{-3/2}
=1-\frac32\varepsilon^2+O(\varepsilon^4).
\]

代入式（12）：

\[
\begin{aligned}
[\cdots]
&=-(1+2\varepsilon+3\varepsilon^2)
-(1-2\varepsilon+3\varepsilon^2)\\
&\quad+2\left(1-\frac32\varepsilon^2\right)
+O(\varepsilon^4)\\
&=-9\varepsilon^2+O(\varepsilon^4).
\end{aligned}
\]

所以首个非零项为

\[
\boxed{
\bm E
\simeq-\frac{9k_eqR_0^2}{l^4}\,\bm e_x
}.
\tag{13}
\]

负号表明电场指向 \(O\)，即沿 \(-x\) 方向。

\textbf{图中初速度与共同构型}

四个初速度必须是绕 \(O\)
顺时针的切向：上方电荷向右，右方电荷向下，下方向左，左方向上。斜指向圆心的箭头会引入径向初速度，改变轨道。由于每个电荷的合力都指向
\(O\)，且力系数相同，四条轨道是彼此转过 \(90^\circ\)
的同一开普勒椭圆；任意同一时刻仍构成正方形。

从近地点 \(R_0\) 到远地点 \(R_1\)，半长轴
\(a=(R_0+R_1)/2\)。周期公式用的是单个电荷的有效中心力常数，不是把四个质量相加。将
\(R_1=R_0\)
代入可还原圆轨道半周期，作为时间式的独立检查。远场一阶、二阶项消失分别来自净电荷为零和偶极矩为零，首项是四极场。

\textbf{力系数与轨道参数逐项核对}

对上方的正电荷，两相邻负电荷距它为 \(\sqrt2r\)，每个吸引力为
\(k_eq^2/(2r^2)\)，向心投影各乘 \(1/\sqrt2\)；对顶的同号电荷距它为
\(2r\)，排斥力为 \(k_eq^2/(4r^2)\)。相减得到

\[F_r=-\frac{k_eq^2}{r^2}\left(\frac1{\sqrt2}-\frac14\right)=-\frac\Gamma{r^2}.\]

有效势 \(U=-\Gamma/r\)。用转向点 \(R_0,R_1\)
的相等总能量消去角动量，可以独立求

\[E=-\frac{\Gamma}{R_0+R_1},\qquad L^2=\frac{2m\Gamma R_0R_1}{R_0+R_1}.\]

这与半长轴 \(a=(R_0+R_1)/2\)、偏心率 \(e=(R_1-R_0)/(R_1+R_0)\)
的开普勒关系一致。时间从一个转向点到另一个是半周期，不能因为空间上只走椭圆的一部分就按弧长比例取时。

四个电荷的相互作用能在求整个系统时每对只计一次；但这里求单个电荷的中心力轨迹，可直接使用上式，不需要再将力系数除以二。能量计数和单粒子运动方程应分别处理。

\newpage

\section{第 6 题}\label{ux7b2c-6-ux9898}

冥王星可视为半径为 \(R\)、质量为 \(M\)
的球体，表面覆盖一层海水，忽略潮汐影响。冥王星绕 \(y\) 轴自转，角速度为
\(\omega\)，万有引力常量为 \(G\)。

\begin{enumerate}
\def\labelenumi{\arabic{enumi}.}
\tightlist
\item
  对海洋表面不在两极的一小质量 \(\Delta m\)，比较其所受万有引力
  \(F_{\text{万}}\) 与支持力 \(N\) 的大小；
\item
  假设海水总量很少且自转不太快，把海洋表面近似为旋转椭球面，其轴截面为椭圆。求椭圆半长轴与半短轴之比；
\item
  求赤道与两极海水深度差 \(\Delta h\)；
\item
  冥王星自转角速度至少多大，赤道处海水才会飞入太空。
\end{enumerate}

\begin{figure}[H]
\centering
\includegraphics[width=0.703\linewidth,height=\textheight,keepaspectratio,alt={图 6-1　冥王星固体表面与旋转海洋等势面。}]{../figures/question-06.pdf}
\caption*{图 6-1　冥王星固体表面与旋转海洋等势面。}
\end{figure}

\subsection{一、题型定位}\label{ux4e00ux9898ux578bux5b9aux4f4d-5}

转动会改变海面等势面，也会改变测得的重力方向。分别处理有效重力、纬度关系和体积约束，避免把海面高度变化与整体零点平移混为一式。

\subsection{二、知识点}\label{ux4e8cux77e5ux8bc6ux70b9-5}

\begin{enumerate}
\def\labelenumi{\arabic{enumi}.}
\tightlist
\item
  旋转质点的向心加速度；
\item
  万有引力、支持力的矢量合成；
\item
  旋转参考系中的离心势；
\item
  小形变近似与椭圆方程；
\item
  赤道失重条件。
\end{enumerate}

\subsection{三、物理图像}\label{ux4e09ux7269ux7406ux56feux50cf-5}

在非极点，海水质点绕 \(y\)
轴做圆周运动。万有引力大体指向球心，支持力沿海面法线向外，两者合力提供向转轴的向心力。因此引力一般大于支持力。

海面稳定时沿面没有切向合力，即海面必须是有效势

\[
\Phi_{\rm eff}
=-\frac{GM}{\sqrt{x^2+y^2}}
-\frac12\omega^2x^2
\]

的等势面。自转使赤道方向鼓起，故半长轴沿 \(x\)，半短轴沿 \(y\)。

\subsection{四、解题入口}\label{ux56dbux89e3ux9898ux5165ux53e3-5}

第（1）问用合力必须指向转轴作定性判断。第（2）问从等势条件出发，利用海面接近球面，把离心项中的半径近似为
\(R\)，整理成椭圆方程。第（3）问令极半径等于固体半径
\(R\)。第（4）问在赤道令支持力降为零。

\subsection{五、解析}\label{ux4e94ux89e3ux6790-5}

\begin{figure}[H]
\centering
\includegraphics[width=0.735\linewidth,height=\textheight,keepaspectratio,alt={图 6-2　旋转参考系中的引力、离心力与海面法向。}]{../figures/solution-06.pdf}
\caption*{图 6-2　旋转参考系中的引力、离心力与海面法向。}
\end{figure}

\textbf{（1）引力与支持力}

质点不在转轴上，随冥王星自转具有指向转轴的向心加速度。若只有万有引力与海面支持力，则

\[
\bm F_{\text{万}}+\bm N
=\Delta m\,\bm a_{\rm c}.
\]

\(\bm F_{\text{万}}\) 指向球心，\(\bm N\)
大体向外。为了留下指向内部的合力，引力大小必须大于支持力。因此

\[
\boxed{F_{\text{万}}>N}.
\tag{1}
\]

对应选择 A。

\textbf{（2）海面椭圆的轴比}

在随冥王星共转的参考系中，稳定海面满足

\[
-\frac{GM}{r}-\frac12\omega^2x^2=C_0,
\qquad
r=\sqrt{x^2+y^2}.
\tag{2}
\]

等价写成

\[
\frac{GM}{r}+\frac12\omega^2x^2=C,
\tag{3}
\]

其中 \(C\) 为常量。因为自转不快，海面只略偏离半径 \(R\)
的球面。对式（3）微分：

\[
-\frac{GM}{r^2}\,\mathrm dr+\omega^2x\,\mathrm dx=0.
\tag{4}
\]

而

\[
r\,\mathrm dr=x\,\mathrm dx+y\,\mathrm dy.
\tag{5}
\]

把式（5）代入式（4）：

\[
-\frac{GM}{r^3}
(x\,\mathrm dx+y\,\mathrm dy)
+\omega^2x\,\mathrm dx=0.
\]

整理：

\[
\left(1-\frac{\omega^2r^3}{GM}\right)x\,\mathrm dx
+y\,\mathrm dy=0.
\tag{6}
\]

小形变近似下，在系数中取 \(r\simeq R\)：

\[
\left(1-\frac{\omega^2R^3}{GM}\right)x\,\mathrm dx
+y\,\mathrm dy=0.
\tag{7}
\]

积分：

\[
\left(1-\frac{\omega^2R^3}{GM}\right)x^2+y^2=C_1.
\tag{8}
\]

与椭圆标准式

\[
\frac{x^2}{a^2}+\frac{y^2}{b^2}=1
\]

比较，有

\[
\frac{a^2}{b^2}
=\frac1{1-\omega^2R^3/(GM)}.
\]

因此

\[
\boxed{
\frac ab
=\frac1{\sqrt{1-\dfrac{\omega^2R^3}{GM}}}
}.
\tag{9}
\]

\textbf{（3）赤道与两极水深差}

海水总量很少，可把极区海面半径近似取为固体半径：

\[
b\simeq R.
\tag{10}
\]

由式（9）

\[
a=\frac{R}{\sqrt{1-\omega^2R^3/(GM)}}.
\]

赤道比极区高出的水深为

\[
\Delta h=a-b.
\]

所以

\[
\boxed{
\Delta h
=R\left[
\frac1{\sqrt{1-\dfrac{\omega^2R^3}{GM}}}-1
\right]
}.
\tag{11}
\]

在小形变精度内，式（11）还可等价写成

\[
\Delta h
\simeq R\left(1-\sqrt{1-\frac{\omega^2R^3}{GM}}\right).
\tag{12}
\]

两式在保留 \(\omega^2\) 的最低阶时相同：

\[
\Delta h\simeq\frac{\omega^2R^4}{2GM}.
\tag{13}
\]

由于本题明确采用``小形变近似''，式（13）是最稳妥的近似结果；式（11）是由近似椭圆方程继续代数求得的表达式。

\textbf{（4）赤道海水飞离条件}

赤道处海水刚好失去支持时 \(N=0\)，万有引力恰好提供向心力：

\[
\frac{GM\Delta m}{R^2}
=\Delta m\,\omega_{\min}^2R.
\]

约去 \(\Delta m\)：

\[
\omega_{\min}^2=\frac{GM}{R^3}.
\]

因此

\[
\boxed{
\omega_{\min}=\sqrt{\frac{GM}{R^3}}
}.
\tag{14}
\]

相应赤道线速度为

\[
v_{\min}=\omega_{\min}R
=\sqrt{\frac{GM}{R}}.
\tag{15}
\]

\textbf{支持力大小的矢量检查}

定义 \(q=\omega^2R^3/(GM)\)，并以余纬 \(\vartheta\)
表示位置。在忽略形变的球面近似中，重力大小
\(F_g=GMm/R^2\)，离心力的大小为
\(qF_g\sin\vartheta\)。支持力与重力、离心力平衡，计算模长得

\[\frac{N^2}{F_g^2}=1-q(2-q)\sin^2\vartheta.\]

所以 \(0<q<1\) 时非极点有
\(N<F_g\)，两极相等。这个结论来自矢量合成，不能仅凭``合力向心''就比较两个不共线力的大小。

形状比 \(a/b\simeq1+q/2\) 是慢转动一阶结果。把同一步近似整理成
\((1-q)^{-1/2}\) 不能把它升级为有限 \(q\)
的精确结果。赤道支持力为零的角速度与逃逸角速度也不同：前者使圆轨道条件成立，后者要求总机械能达到零。

\textbf{等势面方法检查赤道高度差}

在随自转参考系中，有效势为

\[\Phi(r,\vartheta)=-\frac{GM}{r}-\frac12\omega^2r^2\sin^2\vartheta.\]

海面满足 \(\Phi=\text{常量}\)。以极点半径 \(b\) 为基准，赤道半径
\(a=b+\Delta h\)，展开到 \(\omega^2\) 一阶：

\[-\frac{GM}{b}+\frac{GM}{b^2}\Delta h-\frac12\omega^2b^2=-\frac{GM}{b}.\]

从而 \(\Delta h=\omega^2b^4/(2GM)\)，将 \(b\) 替换为未形变半径 \(R\)
只改变更高阶项，正好还原正文高度差。离心势为负号：赤道有效势降低，所以为了保持同一势值，海面需移到重力势较高的位置，即向外升高。

这条路线只检查赤道与极点之差，不确定整颗行星的平均半径；平均半径还需要体积约束。若采用椭球表面，则体积恒定条件
\(a^2b=R^3\) 与一阶 \(a-b\)
共同给出赤道增量和极半径减量。不能把全部高度差都当成赤道相对于原球面的增高。

\newpage

\section{第 7 题}\label{ux7b2c-7-ux9898}

制作一个水火箭。塑料瓶体积为 \(V_0\)，瓶中装有体积 \((1-\alpha)V_0\)
的水，其中
\(0<\alpha<1\)。先向瓶内充入高压气体，高压气体把水在短时间内全部喷出并推动火箭前进。气体视为双原子理想气体，定容摩尔热容
\(C_V=\frac52R\)，大气压为 \(P_0\)。

\begin{enumerate}
\def\labelenumi{\arabic{enumi}.}
\tightlist
\item
  打气机持续提供压强为 \(P_1\)、温度为 \(T_0\)
  的高压气体并缓慢充入塑料瓶。瓶壁很薄且瓶内有水，瓶内气体温度可视为恒定
  \(T_0\)。充到瓶内压强
  \(P_2=\beta P_1\)（\(\beta<1\)）时，求打气机提供的高压气体体积
  \(V_1/V_0\)；
\item
  设单位时间内从打气机打出的体积与压强差满足
  \(\dot V=\gamma\Delta P\)，求充到 \(P_2\) 所需时间；
\item
  发射时气体绝热膨胀并把水在短时间内全部喷出，求喷出水的速度
  \(v_1\)。水密度为 \(\rho_0\)，除水外火箭总质量为 \(M\)；
\item
  取
\end{enumerate}

\[
\alpha=0.75,\quad
\beta=0.20,\quad
\rho_0=1.0\times10^3\,\mathrm{kg/m^3},
\]

\[
M=0.200\,\mathrm{kg},\quad
V_0=2.0\,\mathrm L,\quad
P_1=1.0\times10^6\,\mathrm{Pa},
\]

求第（3）问中水火箭的发射速度 \(v_2\)。

\begin{figure}[H]
\centering
\includegraphics[width=0.678\linewidth,height=\textheight,keepaspectratio,alt={图 7-1　水火箭充气阶段与快速喷射阶段。}]{../figures/question-07.pdf}
\caption*{图 7-1　水火箭充气阶段与快速喷射阶段。}
\end{figure}

\subsection{一、题型定位}\label{ux4e00ux9898ux578bux5b9aux4f4d-6}

先由充气曲线求瓶内初压，再算绝热膨胀释放的能量。速度计算采用原参考的整体两体模型；是否计入外界压强功，会改变可分配的能量。

\subsection{二、知识点}\label{ux4e8cux77e5ux8bc6ux70b9-6}

\begin{enumerate}
\def\labelenumi{\arabic{enumi}.}
\tightlist
\item
  理想气体状态方程；
\item
  以供气状态折算气体体积；
\item
  压差控制的流量微分方程；
\item
  双原子气体绝热指数 \(\kappa=7/5\)；
\item
  绝热关系 \(TV^{\kappa-1}=\text{常量}\)；
\item
  气体内能变化；
\item
  短时喷射过程的动量守恒与能量守恒。
\end{enumerate}

\subsection{三、物理图像}\label{ux4e09ux7269ux7406ux56feux50cf-6}

充气阶段水量不变，瓶内气体体积始终为 \(\alpha V_0\)，温度近似
\(T_0\)，压强从 \(P_0\) 上升到
\(P_2\)。打气机每送入一小体积高压气体，瓶内物质的量增加。

发射阶段极短，气体从 \(\alpha V_0\) 膨胀到
\(V_0\)，水向下喷出，火箭向上运动。在整体两体近似且忽略短时外力冲量时，水与火箭总动量守恒；忽略外界压强功后，气体内能下降才等于两部分动能之和。

\subsection{四、解题入口}\label{ux56dbux89e3ux9898ux5165ux53e3-6}

第（1）问分别写初态、末态和输入气体的物质的量。第（2）问用

\[
\alpha V_0\,\mathrm dP=P_1\,\mathrm dV
\]

把瓶内压强变化与供气体积联系，再代入
\(\mathrm dV/\mathrm dt=\gamma(P_1-P)\)。第（3）问先由绝热关系求末温度，再求气体内能下降；随后联立动量和能量。第（4）问代数值。

\subsection{五、解析}\label{ux4e94ux89e3ux6790-6}

\begin{figure}[H]
\centering
\includegraphics[width=0.696\linewidth,height=\textheight,keepaspectratio,alt={图 7-2　充气压强随时间变化与喷射动量分配。}]{../figures/solution-07.pdf}
\caption*{图 7-2　充气压强随时间变化与喷射动量分配。}
\end{figure}

\textbf{（1）供气体积}

瓶内初态气体体积为 \(\alpha V_0\)，压强为大气压 \(P_0\)，温度为
\(T_0\)。初态物质的量

\[
n_0=\frac{P_0\alpha V_0}{RT_0}.
\tag{1}
\]

充气末态压强为 \(P_2=\beta P_1\)，体积仍为 \(\alpha V_0\)，物质的量

\[
n_2=\frac{P_2\alpha V_0}{RT_0}.
\tag{2}
\]

打气机提供的气体以 \((P_1,T_0)\) 状态折算体积为 \(V_1\)，对应物质的量

\[
n_{\rm in}=\frac{P_1V_1}{RT_0}.
\tag{3}
\]

物质的量守恒：

\[
n_0+n_{\rm in}=n_2.
\]

代入式（1）至式（3）并乘以 \(RT_0\)：

\[
P_0\alpha V_0+P_1V_1=P_2\alpha V_0.
\]

所以

\[
\frac{V_1}{V_0}
=\alpha\frac{P_2-P_0}{P_1}.
\]

再用 \(P_2=\beta P_1\)：

\[
\boxed{
\frac{V_1}{V_0}
=\alpha\left(\beta-\frac{P_0}{P_1}\right)
}.
\tag{4}
\]

\textbf{（2）充气时间}

瓶内温度和体积保持不变，由

\[
P\alpha V_0=nRT_0
\]

微分得

\[
\alpha V_0\,\mathrm dP=RT_0\,\mathrm dn.
\tag{5}
\]

打气机送入体积微元 \(\mathrm dV\) 的高压气体，物质的量

\[
\mathrm dn=\frac{P_1\,\mathrm dV}{RT_0}.
\tag{6}
\]

把式（6）代入式（5）：

\[
\alpha V_0\,\mathrm dP=P_1\,\mathrm dV.
\tag{7}
\]

给定体积流量

\[
\frac{\mathrm dV}{\mathrm dt}
=\gamma(P_1-P).
\tag{8}
\]

由式（7）

\[
\frac{\alpha V_0}{P_1}\frac{\mathrm dP}{\mathrm dt}
=\gamma(P_1-P).
\]

移项：

\[
\frac{\mathrm dP}{P_1-P}
=\frac{\gamma P_1}{\alpha V_0}\,\mathrm dt.
\tag{9}
\]

从 \(P=P_0\) 积分到 \(P=P_2=\beta P_1\)：

\[
\int_{P_0}^{P_2}
\frac{\mathrm dP}{P_1-P}
=\frac{\gamma P_1}{\alpha V_0}\int_0^t\mathrm dt.
\]

左侧积分为

\[
\left[-\ln(P_1-P)\right]_{P_0}^{P_2}
=\ln\frac{P_1-P_0}{P_1-P_2}.
\]

故

\[
t
=\frac{\alpha V_0}{\gamma P_1}
\ln\frac{P_1-P_0}{P_1-P_2}.
\]

代入 \(P_2=\beta P_1\)：

\[
\boxed{
t=\frac{\alpha V_0}{\gamma P_1}
\ln\frac{P_1-P_0}{(1-\beta)P_1}
}.
\tag{10}
\]

\textbf{（3）水的喷出速度}

双原子理想气体

\[
C_V=\frac52R,
\qquad
C_P=C_V+R=\frac72R.
\]

绝热指数为

\[
\kappa=\frac{C_P}{C_V}=\frac75.
\tag{11}
\]

发射前气体温度为 \(T_0\)、体积为 \(\alpha V_0\)；水全部喷出后气体体积为
\(V_0\)，温度记为 \(T_f\)。绝热关系

\[
T_0(\alpha V_0)^{\kappa-1}
=T_fV_0^{\kappa-1}.
\]

所以

\[
T_f=T_0\alpha^{\kappa-1}
=T_0\alpha^{2/5}.
\tag{12}
\]

发射前瓶内气体物质的量由

\[
P_2\alpha V_0=nRT_0
\]

给出：

\[
n=\frac{P_2\alpha V_0}{RT_0}.
\tag{13}
\]

气体内能减少量为

\[
\begin{aligned}
\Delta U
&=nC_V(T_0-T_f)\\
&=\frac{P_2\alpha V_0}{RT_0}
\cdot\frac52R
\cdot T_0(1-\alpha^{2/5})\\
&=\frac52P_2\alpha V_0(1-\alpha^{2/5})\\
&=\frac52\alpha\beta P_1V_0(1-\alpha^{2/5}).
\end{aligned}
\tag{14}
\]

水质量为

\[
m_w=(1-\alpha)V_0\rho_0.
\tag{15}
\]

忽略短时外力冲量，取向上为正。火箭向上速度为 \(v_2\)，水向下速率为
\(v_1\)，动量守恒：

\[
Mv_2-m_wv_1=0.
\]

即

\[
m_wv_1=Mv_2.
\tag{16}
\]

在题设模型中，气体内能减少转化为水和火箭的动能：

\[
\Delta U
=\frac12m_wv_1^2+\frac12Mv_2^2.
\tag{17}
\]

由式（16）

\[
v_2=\frac{m_w}{M}v_1.
\tag{18}
\]

代入式（17）：

\[
\begin{aligned}
\Delta U
&=\frac12m_wv_1^2
+\frac12M\left(\frac{m_w}{M}\right)^2v_1^2\\
&=\frac12v_1^2
\left(m_w+\frac{m_w^2}{M}\right)\\
&=\frac12v_1^2
\frac{m_w(M+m_w)}{M}.
\end{aligned}
\]

所以

\[
v_1^2
=\frac{2\Delta U\,M}{m_w(M+m_w)}.
\tag{19}
\]

代入式（14）、（15）：

\[
\boxed{
v_1=
\sqrt{
\frac{
5(1-\alpha^{2/5})\alpha\beta P_1V_0M
}{
(1-\alpha)V_0\rho_0
\left[M+(1-\alpha)V_0\rho_0\right]
}
}
}.
\tag{20}
\]

\textbf{（4）火箭发射速度}

由式（16）

\[
v_2=\frac{m_w}{M}v_1.
\]

也可直接写成

\[
\boxed{
v_2=
\sqrt{
\frac{
5(1-\alpha^{2/5})\alpha\beta P_1V_0m_w
}{
M(M+m_w)
}
}
}.
\tag{21}
\]

先计算水质量：

\[
\begin{aligned}
m_w
&=(1-0.75)(2.0\times10^{-3})(1.0\times10^3)\\
&=0.500\,\mathrm{kg}.
\end{aligned}
\]

气体内能减少量：

\[
\begin{aligned}
\Delta U
&=\frac52(0.75)(0.20)(10^6)(2.0\times10^{-3})
\left(1-0.75^{2/5}\right)\\
&\approx81.52\,\mathrm J.
\end{aligned}
\]

于是

\[
\begin{aligned}
v_1
&=\sqrt{
\frac{2(81.52)(0.200)}
{0.500(0.500+0.200)}
}\\
&\approx9.65\,\mathrm{m/s},
\end{aligned}
\]

以及

\[
\begin{aligned}
v_2
&=\frac{0.500}{0.200}\times9.65\\
&\approx24.1\,\mathrm{m/s}.
\end{aligned}
\]

所以

\[
\boxed{v_1\approx9.65\,\mathrm{m/s}},
\qquad
\boxed{v_2\approx24.1\,\mathrm{m/s}}.
\tag{22}
\]

\textbf{喷射速度的模型条件与外界压强功}

下面第3、4问沿用原参考解的整体两体模型：所有喷出的水视为速度相同的一个整体，忽略气体质量、重力冲量和外界压强功，并把气体内能减少
\(\Delta U\)
全部分配为水与火箭的动能。这些条件并非实际逐时喷射过程的自动结论。真实流出水具有随时刻变化的速度；仅给初末气体状态，不能唯一给出喷口的瞬时速度历史。

若仍使用整体模型，但计入恒定大气压对膨胀做的功，则可用动能

\[W=\Delta U-P_0(1-\alpha)V_0,\qquad \Delta U=\frac{\beta P_1\alpha V_0}{\kappa-1}(1-\alpha^{\kappa-1}),\quad\kappa=7/5.\]

水质量 \(m_w=\rho_0(1-\alpha)V_0\)，所以

\[v_1=\sqrt{\frac{2WM}{m_w(M+m_w)}},\qquad v_2=\sqrt{\frac{2Wm_w}{M(M+m_w)}}.\]

给定参数下 \(\Delta U\simeq81.52\,\mathrm J\)；忽略外压时
\(v_1\simeq9.65\,\mathrm{m/s}\)、\(v_2\simeq24.13\,\mathrm{m/s}\)。若另取
\(P_0=1.00\times10^5\,\mathrm{Pa}\)，外压功为
\(50.00\,\mathrm J\)，\(W\simeq31.52\,\mathrm J\)，对应约 \(6.00\) 与
\(15.00\,\mathrm{m/s}\)。这一修正相当大，不能称为很小的偏差。原题未给
\(P_0\) 数值，因而后一组是明确增加参数后的比较。

充气可达条件为 \(P_0<\beta P_1<P_1\)；压强曲线从 \(P_0\) 起步，渐近
\(P_1\)。

\textbf{整体速度为何不能代替真实喷口速度}

在整体模型中，水的总动量为 \(-m_wv_1\)，总动能为
\(m_wv_1^2/2\)；动量守恒给 \(Mv_2=m_wv_1\)，能量守恒给两体动能之和为
\(W\)。消去 \(v_1\) 后

\[K_{\text{火箭}}=\frac{m_w}{M+m_w}W,\qquad K_{\text{水}}=\frac M{M+m_w}W.\]

这两个比例只是两体模型的能量分配。实际水在不同时间排出，若出口绝对速度为
\(u(t)\)、质量流率为 \(\dot m(t)\)，则水动量和动能分别是

\[P_{\text{水}}=\int u(t)\dot m(t)dt,\qquad K_{\text{水}}=\frac12\int u(t)^2\dot m(t)dt.\]

一般不能将二者同时写成同一个平均速度的表达，因为
\(\langle u^2\rangle\ne\langle u\rangle^2\)。所以参考解的 \(v_1\)
是整体近似的速度，而非每一滴水离开喷口时都有的速度。

若要进一步建立逐时模型，需要喷口面积、瓶内气体体积与剩余水质量的关系，并用气体压强随体积变化计算相对流出速度，再积分变质量火箭方程。原题没有给这一整套边界条件，本详解因此明确保留原参考的条件性模型，并把大气压功单独比较。

\newpage

\section{第 8 题}\label{ux7b2c-8-ux9898}

有一块平凸透镜，凸面为绕 \(y\) 轴旋转的旋转椭球面，轴截面方程为

\[
\frac{x^2}{a^2}+\frac{y^2}{b^2}=1.
\]

用波长为 \(\lambda\) 的单色光沿 \(-y\) 方向入射。

\begin{enumerate}
\def\labelenumi{\arabic{enumi}.}
\tightlist
\item
  把平凸透镜放在一片刚性玻璃板上，在 \(x\ll a,b\)
  范围内得到牛顿环。求从内到外第 \(n\) 条亮条纹的半径 \(r_n\)；
\item
  把透镜放在表面张力很大的透明液体上。液体表面除与透镜接触的部分外近似为圆锥面，液体与透镜完全不浸润。在
  \(x\ll a,b\) 范围内得到牛顿环。求从内到外第 \(n\) 条亮条纹的外半径
  \(r_n'\)。液体表面张力系数为 \(\sigma\)，透镜质量为
  \(m\)，重力加速度为 \(g\)。已知当 \(|z|\ll1\) 时，
\end{enumerate}

\[
(1+z)^\alpha\simeq1+\alpha z.
\]

\begin{figure}[H]
\centering
\includegraphics[width=0.774\linewidth,height=\textheight,keepaspectratio,alt={图 8-1　平凸透镜在刚性玻璃板和透明液面上的放置。}]{../figures/question-08.pdf}
\caption*{图 8-1　平凸透镜在刚性玻璃板和透明液面上的放置。}
\end{figure}

\subsection{一、题型定位}\label{ux4e00ux9898ux578bux5b9aux4f4d-7}

透镜通过接触圆的表面张力获得支承，接触圆又决定薄膜厚度和干涉环位置。先求接触几何，再逐层检查反射相位与亮环级次。

\subsection{二、知识点}\label{ux4e8cux77e5ux8bc6ux70b9-7}

\begin{enumerate}
\def\labelenumi{\arabic{enumi}.}
\tightlist
\item
  椭圆在顶点附近的二阶展开；
\item
  反射光牛顿环的半波损失；
\item
  亮纹条件 \(2h=(n-\frac12)\lambda\)；
\item
  椭圆顶点曲率半径 \(\rho=a^2/b\)；
\item
  表面张力竖直分量的环线积分；
\item
  二次方程配方与接触圆半径。
\end{enumerate}

\subsection{三、物理图像}\label{ux4e09ux7269ux7406ux56feux50cf-7}

椭球凸面最低点附近的高度差不是 \(x^2/(2R)\) 中的球面半径 \(R\)，而是

\[
h(x)\simeq\frac{b}{2a^2}x^2
=\frac{x^2}{2\rho},
\qquad
\rho=\frac{a^2}{b}.
\]

放在刚性平板上时，薄膜厚度只由椭球面决定。放在液面上时，液面在接触线处与透镜相切，局部近似为斜直线；接触线是与原光轴同心的圆，所有亮环的中心仍在原光轴上。

\subsection{四、解题入口}\label{ux56dbux89e3ux9898ux5165ux53e3-7}

第（1）问先展开椭圆下半支，求空气膜厚，再用反射亮纹条件。第（2）问先由透镜重力与接触线表面张力的竖直分量平衡求小角
\(\theta\)；然后写出接触点处的椭球切线，计算切线与椭球之间的膜厚，整理为完全平方。

\subsection{五、解析}\label{ux4e94ux89e3ux6790-7}

\begin{figure}[H]
\centering
\includegraphics[width=0.774\linewidth,height=\textheight,keepaspectratio,alt={图 8-2　顶点局部膜厚与液面接触线的受力。}]{../figures/solution-08.pdf}
\caption*{图 8-2　顶点局部膜厚与液面接触线的受力。}
\end{figure}

\textbf{（1）刚性玻璃板上的亮环}

椭圆下半支为

\[
y=-b\sqrt{1-\frac{x^2}{a^2}}.
\tag{1}
\]

最低点在 \(x=0\)、\(y=-b\)。距最低点的高度即空气膜厚：

\[
h(x)=y-(-b)
=b\left(1-\sqrt{1-\frac{x^2}{a^2}}\right).
\tag{2}
\]

当 \(x\ll a\) 时，使用

\[
\sqrt{1-z}\simeq1-\frac z2
\]

得

\[
\begin{aligned}
h(x)
&\simeq b\left[
1-\left(1-\frac{x^2}{2a^2}\right)
\right]\\
&=\frac{b}{2a^2}x^2.
\end{aligned}
\tag{3}
\]

反射光中只有一次半波损失，因此从中心向外第 \(n\) 条亮纹满足

\[
2h=\left(n-\frac12\right)\lambda,
\qquad n=1,2,3,\ldots
\tag{4}
\]

将式（3）代入：

\[
\frac{b}{a^2}r_n^2
=\left(n-\frac12\right)\lambda.
\]

所以

\[
\boxed{
r_n
=a\sqrt{\frac{(n-\frac12)\lambda}{b}}
}.
\tag{5}
\]

\textbf{（2）液面上的接触角}

椭圆最低点的曲率半径为

\[
\rho=\frac{a^2}{b}.
\tag{6}
\]

设液面与透镜接触处的切线相对水平倾角为小量 \(\theta\)。接触圆半径为

\[
c\simeq\rho\theta.
\tag{7}
\]

沿接触圆周，单位长度表面张力为 \(\sigma\)，其竖直分量约为
\(\sigma\sin\theta\simeq\sigma\theta\)。圆周总长度为
\(2\pi c\)，故向上的表面张力合力

\[
F_\sigma
=2\pi c\sigma\theta.
\tag{8}
\]

代入 \(c=\rho\theta\)：

\[
F_\sigma=2\pi\sigma\rho\theta^2.
\tag{9}
\]

平衡透镜重力：

\[
mg=2\pi\sigma\rho\theta^2.
\tag{10}
\]

所以

\[
\theta
=\sqrt{\frac{mg}{2\pi\sigma\rho}}
=\sqrt{\frac{mgb}{2\pi\sigma a^2}}.
\tag{11}
\]

接触点横坐标为

\[
x_0=c=\rho\theta=\frac{a^2}{b}\theta.
\tag{12}
\]

\textbf{（3）液面与椭球面之间的膜厚}

椭球下表面在小量近似下可写为

\[
y_e(x)=-b+\frac{b}{2a^2}x^2.
\tag{13}
\]

在 \(x=x_0\) 处斜率为

\[
\left.\frac{\mathrm dy_e}{\mathrm dx}\right|_{x_0}
=\frac{b}{a^2}x_0
=\theta.
\tag{14}
\]

液面近似为过接触点、斜率为 \(\theta\) 的直线：

\[
y_l(x)
=y_e(x_0)+\theta(x-x_0).
\tag{15}
\]

将

\[
y_e(x_0)=-b+\frac{b}{2a^2}x_0^2
\]

代入：

\[
y_l(x)
=-b+\frac{b}{2a^2}x_0^2+\theta x-\theta x_0.
\tag{16}
\]

空气膜厚取两表面纵坐标差的绝对值。用 \(x_0=a^2\theta/b\) 整理：

\[
\begin{aligned}
h(x)
&=y_e(x)-y_l(x)\\
&=\frac{b}{2a^2}x^2-\theta x
+\theta x_0-\frac{b}{2a^2}x_0^2\\
&=\frac{b}{2a^2}x^2-\theta x
+\frac{a^2}{2b}\theta^2.
\end{aligned}
\tag{17}
\]

于是

\[
2h(x)
=\frac{b}{a^2}x^2-2\theta x+\frac{a^2}{b}\theta^2.
\tag{18}
\]

注意右侧恰为完全平方：

\[
2h(x)
=\frac{b}{a^2}
\left(x-\frac{a^2}{b}\theta\right)^2.
\tag{19}
\]

反射亮纹条件仍为

\[
2h=\left(n-\frac12\right)\lambda.
\]

所以

\[
\left|
x-\frac{a^2}{b}\theta
\right|
=a\sqrt{\frac{(n-\frac12)\lambda}{b}}.
\tag{20}
\]

接触圆相对光轴的径向距离为

\[
x_c=\frac{a^2}{b}\theta,
\]

从原光轴量到外侧亮纹的半径为

\[
r_n'
=\frac{a^2}{b}\theta
+a\sqrt{\frac{(n-\frac12)\lambda}{b}}.
\tag{21}
\]

代入式（11）：

\[
\frac{a^2}{b}\theta
=a\sqrt{\frac{mg}{2\pi\sigma b}}.
\]

最终

\[
\boxed{
r_n'
=a\left[
\sqrt{\frac{mg}{2\pi\sigma b}}
+\sqrt{\frac{(n-\frac12)\lambda}{b}}
\right]
}.
\tag{22}
\]

式（22）是以原光轴为中心的亮环外半径。这里 \(x\)
是径向坐标，不是可正负的平移坐标；接触圆外侧 \(x\ge x_c\)
才存在该空气膜，不能把负号根解释为平移后的另一半圆。

\textbf{用俯视图核对``外半径''}

旋转对称性没有破坏，因此不存在某一个方向的中心位移。局部空气膜应写成
\(h(r)=(r-c)^2/(2\rho)\)，并限定在 \(r\ge c\)。反射亮纹给
\(r_n'=c+\sqrt{(n-1/2)\lambda\rho}\)；\(c\)
是接触圆半径。接触圆内的透镜与液体接触区域不具有这一空气膜，不能把配方方程的另一个根直接计作内亮环。图中同时画出轴截面和同心圆俯视示意，便于把径向截面上的``偏移''与平面中心的平移分开。

\textbf{亮环序号、膜厚与相邻环间距}

刚性板上的 \(h=r^2/(2\rho)\)，反射亮纹从 \(n=1\)
开始，\(2h=(n-1/2)\lambda\)。因此

\[r_n^2=(n-1/2)\lambda\rho,\qquad r_{n+1}^2-r_n^2=\lambda\rho.\]

液面接触圆外侧则有

\[r_n'=c+r_n,\qquad (r_{n+1}'-c)^2-(r_n'-c)^2=\lambda\rho.\]

这不是关于偏移圆心的公式，而是把每个同心亮环半径减去接触圆半径后的平方差。若直接计算
\((r_{n+1}')^2-(r_n')^2\)，还会出现
\(2c(r_{n+1}-r_n)\)，所以它不再是常数。这个可观测差异能帮助检查``半径加常数''的物理意义。

局部展开需要 \(r_n'\ll a,b\)，小接触角还需
\(\theta\ll1\)。环数很大时，即使亮纹条件本身仍可用，椭球面二阶展开也可能失效，须使用原椭球高度；不能将渐近公式无条件用于任意序号。轴截面只画一侧会使人误认为中心平移，因此解析图另附俯视同心圆。

\textbf{以接触圆观测反推质量的变式}

若实验中量出接触圆半径 \(c\)，局部切线斜率为
\(\theta=c/\rho\)。在表面张力支承占主导的小角模型中，透镜重量给

\[mg=2\pi c\sigma\theta=\frac{2\pi\sigma c^2}{\rho},\qquad m=\frac{2\pi\sigma c^2}{g\rho}.\]

若只能量出两条亮环外半径 \(r_n',r_{n+1}'\)，可以先利用

\[r_{n+1}'-r_n'=\sqrt{\lambda\rho}\left(\sqrt{n+1/2}-\sqrt{n-1/2}\right)\]

得到 \(\rho\)，再由 \(c=r_n'-\sqrt{(n-1/2)\lambda\rho}\)
得接触圆和质量。这里两环的序号必须已知，否则还有整数级次不确定性。所得
\(c\) 应非负，并满足 \(c/\rho\ll1\)；否则不能继续使用小角力平衡。

若将测量误差传递到质量，固定 \(\rho\) 时有
\(\delta m/m\simeq2\delta c/c\)。当接触圆很小时，即使半径绝对误差很小，质量相对误差也可能很大；这种不稳定来自真实参数关系，而不是牛顿环公式失效。这个变式可以将解题中的几何量与可测结果对应起来。

\textbf{用两条亮环完整回代接触模型}

取明确增加的测量算例：波长 \(\lambda=500\,\mathrm{nm}\)，局部曲率半径
\(\rho=0.100\,\mathrm m\)，接触圆半径
\(c=2.00\,\mathrm{mm}\)。第一、第二条亮环外半径为

\[r_1'=c+\sqrt{\lambda\rho/2}=2.15811\,\mathrm{mm},\qquad r_2'=c+\sqrt{3\lambda\rho/2}=2.27386\,\mathrm{mm}.\]

这两条环的差为
\(0.115747\,\mathrm{mm}\)，与接触圆大小无关。若只知道两条亮环半径及序号，可以先消去
\(c\)：

\[\rho=\frac{(r_2'-r_1')^2}{\lambda(\sqrt{3/2}-\sqrt{1/2})^2}=0.100\,\mathrm m.\]

再代入第一条环，得 \(c=2.00\,\mathrm{mm}\)。把两条环分别代入膜厚式，得到
\(h(r_1')=\lambda/4\)、\(h(r_2')=3\lambda/4\)，均满足带一次半波损失的反射亮纹条件。

若再给表面张力
\(\sigma=0.0700\,\mathrm{N/m}\)、\(g=9.80\,\mathrm{m/s^2}\)，在正文的表面张力支承主导模型内，透镜质量为

\[m=\frac{2\pi\sigma c^2}{g\rho}\simeq1.795\times10^{-6}\,\mathrm{kg}=1.795\,\mathrm{mg}.\]

这里接触斜率
\(c/\rho=0.0200\)，满足小角条件；亮环半径相对透镜几何尺寸也必须足够小，才能用局部二阶展开。质量值是这组附加测量数据在所选支承模型下的反演结果，不能替代原题给定质量，也不能在浮力或其他支承力不可忽略时沿用。这个回代将环序、相位、膜厚和力平衡四条关系放到了同一组数值上。

\textbf{环半径的误差怎样传到质量}

上述反演还可以检查测量精度。记 \(D=r_2'-r_1'\)、\(k=1/(\sqrt3-1)\)，则
\(\rho\propto D^2\)，\(c=r_1'-kD\)。波长固定时，一阶变化为

\[\frac{\delta\rho}{\rho}=\frac{2(\delta r_2'-\delta r_1')}{D},\qquad \delta c=(1+k)\delta r_1'-k\delta r_2'.\]

接触半径和曲率半径由同一对读数求出，误差相互关联，不能把它们当作两次独立测量再相加。固定
\(\sigma,\lambda,g\)，直接对质量式求微分，得到

\[\frac{\delta m}{m}=A\delta r_1'+B\delta r_2',\]

\[A=2\left(\frac{1+k}{c}+\frac1D\right),\qquad B=-2\left(\frac{k}{c}+\frac1D\right).\]

前述数据给
\(A=1.9645\times10^4\,\mathrm{m^{-1}}\)、\(B=-1.8645\times10^4\,\mathrm{m^{-1}}\)。若两环半径各有独立的
\(1.00\,\mu\mathrm m\) 标准不确定度，曲率半径的相对标准不确定度约为
\(2.44\%\)，质量的约为
\(\sqrt{A^2+B^2}\times10^{-6}=2.71\%\)。这些数值只计半径读数误差；波长、表面张力及支承模型的误差须另外处理。

若两个半径同时偏大同一个量
\(\epsilon\)，环差不变，曲率半径不变，而接触半径偏大
\(\epsilon\)，质量的相对偏差为 \(2\epsilon/c\)。取
\(\epsilon=1.00\,\mu\mathrm m\)，该偏差只有
\(0.100\%\)。共同的零点误差与彼此独立的读数误差传播不同，应先说明误差来源。
\end{document}
